% ==========================================================================
% theory.tex -- working notes: why the smoother degrades at p=3 (sometimes)
% and p=4 (always), why additive variants fail, and what actually fixes it.
% % ==========================================================================
% theory.tex -- working notes: why the smoother degrades at p=3 (sometimes)
% and p=4 (always), why additive variants fail, and what actually fixes it.
% % ==========================================================================
% theory.tex -- working notes: why the smoother degrades at p=3 (sometimes)
% and p=4 (always), why additive variants fail, and what actually fixes it.
% % ==========================================================================
% theory.tex -- working notes: why the smoother degrades at p=3 (sometimes)
% and p=4 (always), why additive variants fail, and what actually fixes it.
% \input{theory.tex} at the end of sbm_patch_mg.tex to compile.
% All numerical values reproducible via theory_certificates.py (numpy only).
% ==========================================================================

\section*{Appendix (working notes): definiteness of the SBM boundary terms at the surrogate staircase, and consequences for smoothers}

\subsection*{N.0 Executive summary}

Write $s := \|\mathbf d\|/h$ for the shift length in units of the mesh size.
On a staircase surrogate ($\lambda = 0$) the flat faces have $s \lesssim 1$
and the cells at staircase steps carry \emph{two} extrapolated faces with
$s$ up to $\sqrt 2$ at reentrant pockets. The findings of these notes:

\begin{enumerate}
  \item The Taylor extension of a degree-$p$ polynomial over distance $sh$
        is polynomial extrapolation; its sharp amplification is the
        Chebyshev factor $T_p(1+2s)$, which grows by a factor
        $\rho(s)^2 = \big(\sqrt s + \sqrt{1+s}\big)^2$ \emph{per polynomial
        degree} ($\approx 3.7$ at $s=\tfrac12$, $\approx 7.5$ at
        $s=\sqrt2$). This scaling is confirmed to within a few percent by
        the computed spectra (N.4).
  \item In 1D the method is provably healthy: $p=1$ is unconditionally
        coercive for outward shifts (Lemma N.2), and $p=2$ remains coercive
        up to the absurd shift $s \approx 2.5$ (Lemma N.4). \emph{The
        breakdown is intrinsically multi-dimensional.}
  \item Computable \emph{global indefiniteness certificates} (N.4): the
        symmetric part of the SBM bilinear form is indefinite ---
        on functions supported in a single boundary cell, hence on any mesh
        containing that configuration --- for
        \[
          p = 3:\ \text{only at two-face step cells with } s \to \sqrt 2
          \quad\text{(deep reentrant pockets)},
        \]
        \[
          p = 4:\ \text{from } s \approx 0.5 \text{ at step cells, } s
          \approx 1 \text{ even at flat faces.}
        \]
        This reproduces the empirical pattern exactly: $p\le2$ robust,
        $p=3$ fails \emph{sometimes} (only when the geometry creates deep
        pockets), $p=4$ fails \emph{always} (every staircase contains
        $s \ge 0.5$ step cells). Increasing the penalty $\sigma$ does
        \textbf{not} help (N.4, item v): the offending modes nearly
        annihilate $\mathcal E$, so they do not feel the penalty.
  \item Inward shifts ($\lambda > 0$, true boundary inside the surrogate)
        are far worse: indefinite for all $p \ge 3$ at \emph{every} shift
        length tested, and not repaired by any of the fixes below. For
        $p \ge 3$ the threshold $\lambda$ should be $0$.
  \item Chamfering the staircase with triangles (the proposal of these
        notes' earlier draft) \textbf{rescues $p=3$ but not $p=4$}: it
        removes exactly the deep-pocket configurations where $p=3$ fails,
        but $p=4$ is already indefinite at flat faces with $s \approx 1$,
        which chamfering does not shorten. There is no sliver problem
        (cuts are mesh-defined, all cells uniformly shape-regular), but a
        triangle is thin in the extrapolation direction, which eats part
        of the benefit (N.6).
  \item Truncating the Taylor order to $k=2$ ($k=1$) restores definiteness
        at step cells up to $s \approx 0.75$ (everywhere), at the price of
        reduced boundary consistency $O(h^{k+1})$ --- which defeats the
        purpose of $p=4$ (N.7).
  \item A no-go lemma (N.7b): \emph{any} linear cell-local extension exact
        on $\mathbb Q_p(K)$ coincides with polynomial extrapolation and
        carries the full Chebyshev amplification. Alternatives must import
        extra information: the PDE, the boundary data $g$, or a larger
        support. The PDE-substituted extension (all normal derivatives of
        order $\ge 2$ traded for tangential derivatives and $f$-data)
        keeps full consistency order, gains roughly one polynomial degree
        of definiteness (certified in N.7b), and combined with the chamfer
        (which caps $s$ at $\approx 0.7$ and removes two-face cells) is
        the predicted robust \emph{full-order} $p=4$ configuration.
\end{enumerate}

\subsection*{N.1 Setting}

Model: cell $K = (0,h)^d$, $\mathbb Q_p$ Lagrange polynomials, SBM form
restricted to $K$ and its surrogate faces $F$:
\begin{equation}
  a_K(u,v) = \int_K \nabla u\cdot\nabla v
  - \int_F (\partial_n u)\, v
  - \int_F (\partial_n v)\, \mathcal E u
  + \sigma \int_F \mathcal E u\, \mathcal E v,
  \qquad \sigma = \gamma\, p^2/h,
  \label{eq:cellform}
\end{equation}
with $\mathcal E u(\mathbf x) = \sum_{j=0}^{k} \tfrac1{j!} (\mathbf
d\cdot\nabla)^j u(\mathbf x)$. Throughout, $s = \|\mathbf d\|/h$ and we set
$h=1$ in computations; $\gamma = 4$ unless stated (results are insensitive
to $\gamma$, see N.4).

Two mechanisms could plausibly cause the observed breakdown: (i) elliptic
corner singularities of the surrogate domain (reentrant staircase corners,
interior angle $3\pi/2$, Kondrat'ev exponent $r^{2/3}$), and (ii) loss of
definiteness of the discrete boundary terms. The results below show that
(ii) is dominant and quantitatively predictive; (i) is discussed briefly in
N.8.

\subsection*{N.2 Sharp extrapolation amplification}

\paragraph{Lemma N.1 (extension = extrapolation; sharp growth).}
\emph{Let $q$ be a polynomial of degree $p$ on $[0,h]$ and let the shift be
axis-aligned. For $k \ge p$ the Taylor extension is exact point evaluation,
$\mathcal E q(0) = q(-\delta)$, and}
\[
  \sup\Big\{ |q(-\delta)| \,:\, \|q\|_{L^\infty(0,h)} \le 1 \Big\}
  \;=\; T_p\!\big(1 + 2s\big), \qquad s = \delta/h,
\]
\emph{attained by the Chebyshev polynomial. Moreover, with
$\rho(s) := \sqrt s + \sqrt{1+s}$,}
\[
  T_p(1+2s) \;=\; \tfrac12\Big( \rho(s)^{2p} + \rho(s)^{-2p} \Big).
\]

\emph{Proof.} Exactness: the Taylor polynomial of order $k\ge p$ of a
degree-$p$ polynomial is the polynomial itself. Extremality of $T_p$
outside $[-1,1]$ is classical (Chebyshev; see Rivlin). The identity follows
from $T_p(\cosh\theta) = \cosh(p\theta)$ with
$1+2s = \cosh\big(2\,\mathrm{arcsinh}\sqrt s\big)$ and
$e^{\mathrm{arcsinh}\sqrt s} = \sqrt s + \sqrt{1+s}$. $\square$

The quantity $\rho(s)^2$ is the \emph{amplification per polynomial degree}:

\begin{center}
\begin{tabular}{lccccc}
\hline
 & $s=\tfrac14$ & $s=\tfrac12$ & $s=\tfrac34$ & $s=1$ & $s=\sqrt2$ \\
\hline
$\rho(s)^2$ (per degree)      & $2.6$ & $3.7$ & $4.8$ & $5.8$ & $7.5$ \\
$T_1(1+2s)$ & $1.5$ & $2$ & $2.5$ & $3$ & $3.8$ \\
$T_2(1+2s)$ & $3.5$ & $7$ & $11.5$ & $17$ & $28$ \\
$T_3(1+2s)$ & $9$ & $26$ & $55$ & $99$ & $214$ \\
$T_4(1+2s)$ & $24.5$ & $97$ & $265$ & $577$ & $1611$ \\
\hline
\end{tabular}
\end{center}

The penalty and consistency terms in \eqref{eq:cellform} are bilinear in
$\mathcal E$, so the boundary block scales like $\sigma\, T_p^2$ relative to
the $O(1)$ (in $h^{d-2}$ units) volume stiffness. This is confirmed by the
computed largest eigenvalues $\lambda_{\max}$ of the symmetric part: e.g.\
for the flat configuration at $p=4$, the measured ratios
$\lambda_{\max}(s{=}1)/\lambda_{\max}(s{=}\tfrac12) = 34$ and
$\lambda_{\max}(s{=}\sqrt2)/\lambda_{\max}(s{=}1) = 7.6$ match the predicted
$\big(T_4(3)/T_4(2)\big)^2 = 35$ and $\big(T_4(3.83)/T_4(3)\big)^2 = 7.8$.
At $p=4$, $s=1$ the local block condition number already exceeds $10^7$.
Conditioning alone, however, is not the failure mode --- definiteness is,
as the next two subsections show.

\subsection*{N.3 Exact one-dimensional results}

On $K=(0,h)$ with surrogate face at $x=0$ and true boundary at $-\delta$
(outward shift), $\mathcal E v(0) = v(-\delta)$, form \eqref{eq:cellform}.

\paragraph{Lemma N.2 ($p=1$, closed form).} \emph{For $v = a + bx$:}
\[
  a_K(v,v) = b^2 (h+\delta) + 2bw + \sigma w^2, \qquad w := v(-\delta),
\]
\emph{which is positive definite iff $\sigma (h + \delta) > 1$. For the
inward shift ($\Gamma$ at $+\delta$ inside the cell) the same computation
gives $b^2(h-\delta) + 2bw + \sigma w^2$, positive definite iff
$\sigma (h - \delta) > 1$.}

Consequences: $p=1$ with outward shifts is unconditionally coercive (any
$\sigma \ge 1/h$, any $\delta > 0$) --- matching the total robustness of
$p=1$ in practice. Inward shifts lose coercivity as $\delta \to h$; with
$\sigma = 4/h$ the threshold is $\delta = 3h/4$, reproduced to machine
precision by the 2D code (a useful cross-validation of both).

\paragraph{Lemma N.3 (kernel escape).} \emph{On the subspace
$\{v : \mathcal E v = 0 \text{ on } F\}$ the penalty and one flux term
vanish identically and \eqref{eq:cellform} reduces to the
\emph{unpenalized half-Nitsche} form}
\[
  F(v) = \|\nabla v\|_{L^2(K)}^2 - \int_F (\partial_n v)\, v .
\]
Definiteness on this kernel is a sign-correlation property between $v$ and
$\partial_n v$ on $F$, forced by $v$ vanishing at the shifted points; the
penalty is powerless there, no matter how large. This is the structural
reason for the $\gamma$-insensitivity observed in N.4.

\paragraph{Lemma N.4 ($p=2$, 1D, explicit threshold).} \emph{For $v$
quadratic with $v(-\delta) = 0$ (kernel of $\mathcal E$), writing
$c = v'(0) $-parametrization $v = (x+\delta)(a+bx)$, $h=1$:}
\[
  F(v) = (1+s)\,c^2 + (2 - s^2)\,cb + \tfrac43\, b^2, \qquad c = a + bs,
\]
\emph{positive definite iff $(2-s^2)^2 < \tfrac{16}{3}(1+s)$, i.e.\ for all
$s \lesssim 2.5$.}

So in 1D even $p=2$ survives shifts of two and a half cells; the observed
2D breakdown at $s \le \sqrt2$ cannot be a 1D phenomenon. What 1D misses:
at a staircase step, \emph{two} faces with orthogonal shift directions load
the same cell polynomial, and the sign-correlation of Lemma N.3 cannot be
satisfied for both directions simultaneously once $p$ gives the polynomial
enough freedom to oscillate between the cell and the shifted evaluation
points.

\subsection*{N.4 Numerical definiteness certificates in 2D}

\paragraph{Certificate construction.} Let $K$ be a boundary cell in one of
the configurations below and let $v \in \mathbb Q_p(K)$ vanish on the edges
of $K$ shared with active neighbors. Extending $v$ by zero yields a
conforming global function whose global energy equals $a_K(v,v)$: volume
terms of neighbors vanish, and all other surrogate-face terms vanish
because $\mathcal E$ on a face of cell $K'$ uses only the polynomial of
$K'$. Hence
\[
  \lambda_{\min}\Big( \tfrac12 (A_K + A_K^T)\big|_{S} \Big) < 0
  \;\Longrightarrow\;
  \text{the \emph{global} symmetric part is indefinite}
\]
on every mesh containing the configuration, where $S$ is the zero-extension
subspace. Configurations ($h=1$, exact extrapolation $=$ Taylor with $k=p$
for the axis-aligned shifts):
\emph{flat} (one face $\{y{=}1\}$, $\mathbf d = (0,s)$; $S$: $v=0$ on the
three interior edges), \emph{step} (two faces $\{y{=}1\},\{x{=}1\}$,
$\mathbf d = (0,s),(s,0)$; the staircase step cell; $S$: $v=0$ on two
interior edges), \emph{step-diag} (as step but both shifts
$s(1,1)/\sqrt2$, modeling pocket-directed closest-point projections), and
\emph{chamfer} variants in N.6.

\paragraph{Results ($\gamma = 4$).} Minimal eigenvalue of the symmetric
part on $S$ (negative $=$ certified global indefiniteness; the volume
stiffness is $O(1)$ in these units, so magnitudes are meaningful):

\begin{center}
\begin{tabular}{llcccc}
\hline
config & $p$ & $s=0.5$ & $s=0.75$ & $s=1$ & $s=\sqrt2$ \\
\hline
flat  & 3 & $+1.22$ & $+0.91$ & $+0.75$ & $+0.61$ \\
flat  & 4 & $+0.28$ & $+0.08$ & $\mathbf{-0.03}$ & $\mathbf{-0.13}$ \\
step  & 3 & $+0.27$ & $+0.09$ & $+0.02$ & $\mathbf{-0.05}$ \\
step  & 4 & $\mathbf{-0.04}$ & $\mathbf{-0.13}$ & $\mathbf{-0.18}$ & $\mathbf{-0.24}$ \\
\hline
\end{tabular}
\end{center}

Findings:
\begin{enumerate}
  \item[(i)] $p \le 2$: positive in every configuration and shift tested
        (not shown; margins $O(1)$).
  \item[(ii)] $p=3$: indefinite \emph{only} at the two-face step cell with
        $s \to \sqrt2$, i.e.\ at deep reentrant pockets. Whether a given
        geometry produces such pockets depends on how $\Gamma$ cuts the
        grid --- hence ``$p=3$ sometimes breaks''.
  \item[(iii)] $p=4$: indefinite from $s \approx 0.5$ at step cells and
        from $s \approx 1$ at flat faces. Every staircase surrogate of a
        generic $\Gamma$ contains such cells --- hence ``$p=4$ always
        breaks''.
  \item[(iv)] The nonsymmetric spectra show only small imaginary parts for
        outward shifts: for $\lambda = 0$ the pathology is
        \emph{indefiniteness of the symmetric part}, not complex
        eigenvalues (those belong to inward shifts, consistent with the 1D
        study in~\cite{wichrowski2025geometric}).
  \item[(v)] Penalty insensitivity (Lemma N.3 in action), step cell,
        $s=1$: for $p=4$, $\lambda_{\min}|_S = -0.191, -0.184, -0.178,
        -0.177$ at $\gamma = 2, 4, 16, 64$. No penalty tuning can repair
        this.
  \item[(vi)] Inward shifts (flat face, $\mathbf d = (0,-s)$, corrected
        certificate): $p=2$ positive at all $s$; $p=3$ and $p=4$
        \emph{negative at every} $s \in \{0.25, 0.5, 0.75\}$ (e.g.\ $-0.63$
        and $-1.13$ at $s=0.25$). Intersected-cell surrogates
        ($\lambda>0$) are therefore incompatible with $p \ge 3$
        independently of corner effects.
\end{enumerate}

\subsection*{N.5 Consequences for smoothers}

Let $v_\star$ be a certified negative mode ($a(v_\star,v_\star) < 0$,
supported in one boundary cell).

\begin{itemize}
  \item \textbf{All energy-based convergence theory is void.} Both the
        Schwarz theory for additive/multiplicative smoothers and the
        field-of-values GMRES bounds (Eisenstat--Elman--Schultz,
        Cai--Widlund) presuppose a definite symmetric part. On staircase
        SBM with $p \ge 3$ this hypothesis \emph{fails as a matter of
        fact}, not of proof technique.
  \item \textbf{Pointwise and cell-wise additive smoothers.} For
        (block-)Jacobi-type methods the local blocks themselves become
        indefinite (the single-cell Neumann blocks go negative one degree
        earlier than the certificates, already $p=3$, $s \ge 0.75$).
        Damped Richardson with an indefinite block-diagonal preconditioner
        has no admissible damping: the modes with negative Rayleigh
        quotient are amplified for every $\omega > 0$ of either sign
        globally. This is the cleanest available explanation for the
        observed \emph{divergence} (not merely stagnation) of the
        cell-wise additive DG+SBM smoother.
  \item \textbf{Why multiplicative full-residual patches still work at
        $p \le 3$.} The negative mode $v_\star$ is supported in one cell,
        hence entirely inside at least one vertex patch (the full-residual
        construction guarantees its whole residual row is visible there).
        An \emph{exact} local solve needs only invertibility, not
        definiteness, of $A_i$; visiting that patch removes $v_\star$'s
        error component exactly, and the sequential sweep propagates the
        repaired residual outward. In the additive version the $2^d$
        overlapping patches containing $v_\star$ each compute a correction
        against an indefinite, non-normal block and their \emph{sum} is
        applied blindly; there is no self-correction within the sweep, and
        the usual multiplicity damping $\omega \le 1/N_c$ cannot produce
        contraction where the field of values crosses zero.
  \item \textbf{Why $p=4$ defeats even the multiplicative smoother.} The
        indefinite modes are no longer isolated: at $p=4$ they exist at
        \emph{every} boundary cell with $s \gtrsim 0.5$--$1$, they couple
        to neighboring patches through off-block entries of size
        $O(\sigma T_p^2)$ ($\ge 10^6$ at $s = 1$), and the error injected
        across a patch boundary is re-amplified faster than the sweep
        damps it. Locality of the repair mechanism is lost.
\end{itemize}

\subsection*{N.6 Chamfer (mesh-defined triangular cuts): rescues $p=3$, not $p=4$}

The chamfer splits staircase corner quads along the vertex-to-vertex
diagonal. The cut position is fixed by the mesh: every cut cell is exactly
half a quad, uniformly shape-regular --- \emph{there is no small-cut/sliver
problem and no ghost penalty is needed}, in contrast to
$\Gamma$-interpolating CutFEM cuts. Reentrant $3\pi/2$ corners become pairs
of $5\pi/4$ corners (regularity exponent $2/3 \to 4/5$), and the worst
shift drops from $s \approx \sqrt2$ to $s \approx 0.7$ with a single,
near-normal face replacing the two-face loading.

Certificates for the chamfered cell (hypotenuse face, shift
$s(1,1)/\sqrt2$), two variants: $\mathbb P_p$ on the triangle (remeshed)
and $\mathbb Q_p$ on the parent quad integrated over the triangle
(cut-quadrature):

\begin{center}
\begin{tabular}{llccc}
\hline
variant & $p$ & $s=0.5$ & $s=0.707$ & $s=1$ \\
\hline
$\mathbb P_p$ triangle & 3 & $+0.91$ & $+0.62$ & $+0.42$ \\
$\mathbb P_p$ triangle & 4 & $\mathbf{-0.15}$ & $\mathbf{-0.34}$ & $\mathbf{-0.48}$ \\
$\mathbb Q_p$ cut      & 3 & $\mathbf{-0.01}$ & $\approx\mathbf{-0.03}$ & $\mathbf{-0.02}$ \\
$\mathbb Q_p$ cut      & 4 & $\mathbf{-0.08}$ & $\mathbf{-0.07}$ & $\mathbf{-0.06}$ \\
\hline
\end{tabular}
\end{center}

Reading:
\begin{itemize}
  \item The \emph{remeshed} chamfer ($\mathbb P_p$ triangles) is positive
        for $p=3$ at all relevant shifts. Since the chamfer also caps $s$
        at $\approx 0.7$, it removes exactly the configuration in which
        $p=3$ was certified indefinite. \textbf{Prediction: chamfering
        makes $p=3$ uniformly robust.}
  \item It does \emph{not} rescue $p=4$: the triangle block is indefinite
        from $s = 0.5$. Two reasons: (a) a triangle is thin in the
        extrapolation direction (extent $h/\sqrt2$ normal to the
        hypotenuse, tapering), so the effective Chebyshev argument is
        larger than $s$ suggests; (b) independently, $p=4$ is already
        indefinite at \emph{flat} faces with $s \approx 1$, which the
        chamfer does not touch.
  \item The cut-quadrature variant is weaker (marginally indefinite even
        at $p=3$): the polynomial is energy-controlled only on the
        half-cell while being extrapolated from its edge. If chamfering is
        implemented, the triangles should carry genuine $\mathbb P_p$
        spaces (or the parent-cell variant needs a stabilization on the
        inactive half). This is an implementation-relevant distinction the
        earlier draft of these notes missed.
\end{itemize}

\subsection*{N.7 The effective fix: truncated extrapolation order}

Lemma N.1 identifies the amplification driver as the extrapolation
\emph{order}: truncating the Taylor extension at order $k < p$ caps the
growth at degree-$k$ extrapolation ($T_k$-like) while retaining
$O((sh)^{k+1})$ boundary consistency. Certificates for the step cell
(worst outward configuration), truncation order $k$:

\begin{center}
\begin{tabular}{llcccc}
\hline
config & $p$ & $k=p$ & $k=3$ & $k=2$ & $k=1$ \\
\hline
step, $s=0.75$   & 4 & $-0.13$ & $-0.10$ & $\mathbf{+0.01}$ & $+0.41$ \\
step, $s=1$      & 4 & $-0.18$ & $-0.14$ & $-0.03$ & $+0.35$ \\
step, $s=\sqrt2$ & 4 & $-0.24$ & $-0.18$ & $-0.06$ & $+0.30$ \\
step, $s=\sqrt2$ & 3 & $-0.05$ & --      & $\mathbf{+0.04}$ & $+0.50$ \\
flat, $s=1$      & 4 & $-0.03$ & $+0.06$ & $+0.34$ & $+0.98$ \\
\hline
\end{tabular}
\end{center}

\begin{itemize}
  \item $k=1$ restores definiteness \emph{everywhere} with $O(1)$ margins
        --- a bulletproof but low-order option (locally $O(h^2)$
        consistency; acceptable if applied only at the few worst faces).
  \item $k=2$ restores definiteness for $p=4$ up to $s \approx 0.75$ and
        for $p=3$ everywhere, with $O(h^3)$ boundary consistency.
  \item Honest objection: with a staircase or chamfered surrogate,
        $\delta = O(h)$, so $k<p$ caps the \emph{global} boundary accuracy
        at $O(h^{k+1})$ and largely defeats the purpose of $p=4$. (The
        truncation is order-neutral only where $\delta = O(h^2)$, i.e.\
        with a $\Gamma$-interpolating cut.) A shift-adaptive rule
        ($k=p$ where $s \le 0.4$, smaller $k$ beyond) confines the loss to
        the worst faces, but the principled resolution is to change the
        extension, not its order --- see N.7b.
  \item For inward shifts truncation helps at moderate $s$ but $p \ge 3$
        remains indefinite at small inward $s$ regardless of $k$ ---
        reinforcing finding (vi): use $\lambda = 0$ for $p \ge 3$.
\end{itemize}

\subsection*{N.7b Alternative extension operators: a no-go lemma and three ways around it}

\paragraph{Lemma N.5 (no free lunch for cell-local extensions).}
\emph{Let $\mathcal L : \mathbb Q_p(K) \to \mathbb R$ be linear, computed
from data on $K$ only, and exact on polynomials,
$\mathcal L q = q(\tilde{\mathbf x} + \mathbf d)$ for all $q \in
\mathbb Q_p(K)$ (as required for $O(h^{p+1})$ consistency from cell-local
data). Then $\mathcal L$ \emph{is} the point-evaluation extrapolation on
$\mathbb Q_p(K)$, and its norm is the Chebyshev factor of Lemma N.1.}
\emph{Proof:} two linear functionals agreeing on the whole space are
equal. $\square$

Smoothness of the extension is a red herring: on the discrete space every
exact extension coincides with the one already in use. Any genuine
alternative must therefore use information beyond the single cell's
polynomial. Exactly three sources exist, and each gives a known method
family:

\paragraph{(A) Use the PDE: substituted normal derivatives (recommended).}
Since $-\Delta u = f$, all normal derivatives of order $\ge 2$ can be
eliminated recursively,
\[
  \partial_n^2 u = -f - \Delta_{\mathrm{tan}} u, \qquad
  \partial_n^3 u = -\partial_n f - \Delta_{\mathrm{tan}} \partial_n u,
  \quad \dots
\]
so the degree-$k$ Taylor extension becomes
\[
  \mathcal E_{\mathrm{pde}} u
  = u + \delta\, \partial_n u
  + \sum_{j=2}^{k} \tfrac{\delta^j}{j!}
    \big( \text{tangential derivatives of } u,\ \partial_n u \big)
  + \underbrace{\big( f\text{-terms} \big)}_{\text{data, rhs}} ,
\]
with \emph{every $u$-dependent term evaluated on the face}: no
extrapolation beyond the first normal derivative. Crucially,
$\mathcal E_{\mathrm{pde}}$ is exact on solutions of the PDE --- which is
all that consistency requires --- so \textbf{the full order $O(\delta^{k+1})$
is retained} while the Chebyshev growth is traded for tangential Markov
factors $\big(c\,p^2 \delta/h\big)^{j}$, small for moderate $s$.
Certificates (same setup as N.4, $f$-part irrelevant for the spectrum):

\begin{center}
\begin{tabular}{llcccc}
\hline
config & $p$ & $s=0.5$ & $s=0.75$ & $s=1$ & $s=\sqrt2$ \\
\hline
flat, $\mathcal E_{\mathrm{pde}}$ & 3 & $+2.87$ & $+3.36$ & $+3.84$ & $+4.25$ \\
flat, $\mathcal E_{\mathrm{pde}}$ & 4 & $+1.34$ & $+0.80$ & $-0.01$ & $-0.71$ \\
step, $\mathcal E_{\mathrm{pde}}$ & 3 & $+1.20$ & $+0.90$ & $+0.49$ & $-0.10$ \\
step, $\mathcal E_{\mathrm{pde}}$ & 4 & $+0.12$ & $-0.34$ & $-0.38$ & $-0.33$ \\
\hline
\end{tabular}
\end{center}

Compared with plain Taylor (N.4): roughly one polynomial degree of
definiteness is gained at full consistency order ($p=4$ flat: $+0.08 \to
+0.80$ at $s = 0.75$; step: indefinite from $s=0.5 \to$ definite at
$s=0.5$). At $s \gtrsim 1$ the tangential factors take over and
indefiniteness returns --- so $\mathcal E_{\mathrm{pde}}$ alone does not
cure the staircase, but \textbf{chamfer ($s \le 0.7$, single near-normal
face) $+$ $\mathcal E_{\mathrm{pde}}$ is certified-definite in every
configuration that then occurs, at full order.} Implementation cost:
tangential/mixed derivatives of the shape functions on the face (already
available), plus $f$ and its normal-derivative data near the boundary in
the right-hand side.

\paragraph{(B) Use the boundary data: interpolatory (ghost/ROD-type)
conditions.} Taylor never uses the one thing that is known exactly: $g$ at
the projected point $\mathbf x \in \Gamma$. Writing the boundary condition
via a 1D interpolant along the shift line through $p{+}1$ points
\emph{inside} $\tilde\Omega$ plus the pinned value $g(\mathbf x)$, the
surrogate-face value $u(\tilde{\mathbf x})$ is expressed by
\emph{interpolation} (the evaluation point lies between the samples and
$\mathbf x$), so the $u$-dependent weights are bounded by a Lebesgue
constant instead of $T_p$; full order is retained. This is the
finite-element analogue of ghost-point methods (Gibou--Fedkiw) and
reconstructed off-site data (ROD). Cost: the stencil crosses 1--2 interior
cells, so the sparsity pattern widens; the samples stay inside a vertex
patch, so the full-residual construction survives.

\paragraph{(C) Use a larger support: aggregated extension.} Extrapolate
from a least-squares polynomial fit over an interior patch of diameter
$L = 2h$--$3h$ (AgFEM-flavored). The amplification drops to
$T_p(1 + 2\delta/L)$; by the certificate tables, an effective shift
$\delta/L \le 0.4$ is safe even at $p=4$, so $L \ge 3h$ suffices for the
worst staircase pockets. Cost: same widened coupling as (B), plus the fit;
conceptually the bluntest of the three.

Ranking for this project: (A) is the least invasive (face-local, no
sparsity change, no new data structures) and combines multiplicatively
with the chamfer; (B) is the most robust in principle (bounded weights
regardless of geometry) and is the natural fallback if (A)$+$chamfer still
misbehaves at $p \ge 5$ or in 3D, where edge/vertex pockets give $s$ up to
$\sqrt3$.

\subsection*{N.8 The regularity mechanism (demoted to second order)}

Reentrant staircase corners carry Kondrat'ev singularities
$r^{2/3}\sin(2\theta/3)$; they are recreated at new positions on every
level (non-nested active sets), so the associated error is $O(h)$-wavelength
by construction and falls entirely on the fine-level smoother, and
degree-$p$ approximation saturates at $O(h^{2/3})$ locally. This degrades
constants and localizes the hard error components at the boundary --- it
explains \emph{where} the smoother must work hardest, but not the observed
$p$-thresholds: the certificates of N.4 do, quantitatively and with the
right $\lambda$-, $s$-, and configuration-dependence. If the exact SBM
extrapolation were used ($k$ exact, no indefiniteness), the corner
singularities alone would be handled by the existing patch smoother, as
they are for $p \le 2$.

\subsection*{N.9 What is proven, what is certified, what is conjectured}

\emph{Proven} (Lemmas N.1--N.4): sharp Chebyshev amplification with the
per-degree factor $\rho(s)^2$; unconditional 1D $p=1$ coercivity (outward)
and the inward threshold $\sigma(h-\delta) > 1$; the kernel-escape
mechanism and penalty-insensitivity; 1D $p=2$ coercivity to $s\approx2.5$.

\emph{Certified numerically} (exact statements about finite matrices,
N.4/N.6/N.7): global indefiniteness of the symmetric part for the listed
$(p, s, \text{config})$; its penalty-insensitivity; positivity of the
chamfered $\mathbb P_3$ and truncated-$k$ variants.

\emph{Proven} (Lemma N.5): no cell-local exact extension can beat the
Chebyshev amplification; alternatives must use the PDE, the data $g$, or a
larger support (N.7b).

\emph{Conjectured / to validate in the solver}: (a) chamfer $\Rightarrow$
$p=3$ uniformly robust; (b) chamfer $+$ PDE-substituted extension
$\mathcal E_{\mathrm{pde}}$ $\Rightarrow$ $p=4$ robust \emph{at full
order} (the leading recommendation); (c) shift-adaptive $k$-truncation
alone (no chamfer) already recovers $p=4$ convergence with the existing
smoother at reduced boundary accuracy --- still the cheapest diagnostic
experiment: a ten-line change that directly tests the whole theory;
(d) restoring definiteness restores the additive patch smoother to
usefulness (with standard multiplicity damping), enabling the parallel
variant.

% --- references to add if these notes graduate into the paper:
% Rivlin, "Chebyshev Polynomials" (extremal growth outside [-1,1]).
% Grisvard, "Elliptic Problems in Nonsmooth Domains" (corner asymptotics).
% Atallah, Canuto, Scovazzi (SBM analysis, d/h conditions).
% Cai & Widlund SINUM 1992/93; Eisenstat, Elman, Schultz 1983 (nonsymmetric bounds).
% Burman et al. (CutFEM/ghost penalty, for the contrast drawn in N.6).
 at the end of sbm_patch_mg.tex to compile.
% All numerical values reproducible via theory_certificates.py (numpy only).
% ==========================================================================

\section*{Appendix (working notes): definiteness of the SBM boundary terms at the surrogate staircase, and consequences for smoothers}

\subsection*{N.0 Executive summary}

Write $s := \|\mathbf d\|/h$ for the shift length in units of the mesh size.
On a staircase surrogate ($\lambda = 0$) the flat faces have $s \lesssim 1$
and the cells at staircase steps carry \emph{two} extrapolated faces with
$s$ up to $\sqrt 2$ at reentrant pockets. The findings of these notes:

\begin{enumerate}
  \item The Taylor extension of a degree-$p$ polynomial over distance $sh$
        is polynomial extrapolation; its sharp amplification is the
        Chebyshev factor $T_p(1+2s)$, which grows by a factor
        $\rho(s)^2 = \big(\sqrt s + \sqrt{1+s}\big)^2$ \emph{per polynomial
        degree} ($\approx 3.7$ at $s=\tfrac12$, $\approx 7.5$ at
        $s=\sqrt2$). This scaling is confirmed to within a few percent by
        the computed spectra (N.4).
  \item In 1D the method is provably healthy: $p=1$ is unconditionally
        coercive for outward shifts (Lemma N.2), and $p=2$ remains coercive
        up to the absurd shift $s \approx 2.5$ (Lemma N.4). \emph{The
        breakdown is intrinsically multi-dimensional.}
  \item Computable \emph{global indefiniteness certificates} (N.4): the
        symmetric part of the SBM bilinear form is indefinite ---
        on functions supported in a single boundary cell, hence on any mesh
        containing that configuration --- for
        \[
          p = 3:\ \text{only at two-face step cells with } s \to \sqrt 2
          \quad\text{(deep reentrant pockets)},
        \]
        \[
          p = 4:\ \text{from } s \approx 0.5 \text{ at step cells, } s
          \approx 1 \text{ even at flat faces.}
        \]
        This reproduces the empirical pattern exactly: $p\le2$ robust,
        $p=3$ fails \emph{sometimes} (only when the geometry creates deep
        pockets), $p=4$ fails \emph{always} (every staircase contains
        $s \ge 0.5$ step cells). Increasing the penalty $\sigma$ does
        \textbf{not} help (N.4, item v): the offending modes nearly
        annihilate $\mathcal E$, so they do not feel the penalty.
  \item Inward shifts ($\lambda > 0$, true boundary inside the surrogate)
        are far worse: indefinite for all $p \ge 3$ at \emph{every} shift
        length tested, and not repaired by any of the fixes below. For
        $p \ge 3$ the threshold $\lambda$ should be $0$.
  \item Chamfering the staircase with triangles (the proposal of these
        notes' earlier draft) \textbf{rescues $p=3$ but not $p=4$}: it
        removes exactly the deep-pocket configurations where $p=3$ fails,
        but $p=4$ is already indefinite at flat faces with $s \approx 1$,
        which chamfering does not shorten. There is no sliver problem
        (cuts are mesh-defined, all cells uniformly shape-regular), but a
        triangle is thin in the extrapolation direction, which eats part
        of the benefit (N.6).
  \item Truncating the Taylor order to $k=2$ ($k=1$) restores definiteness
        at step cells up to $s \approx 0.75$ (everywhere), at the price of
        reduced boundary consistency $O(h^{k+1})$ --- which defeats the
        purpose of $p=4$ (N.7).
  \item A no-go lemma (N.7b): \emph{any} linear cell-local extension exact
        on $\mathbb Q_p(K)$ coincides with polynomial extrapolation and
        carries the full Chebyshev amplification. Alternatives must import
        extra information: the PDE, the boundary data $g$, or a larger
        support. The PDE-substituted extension (all normal derivatives of
        order $\ge 2$ traded for tangential derivatives and $f$-data)
        keeps full consistency order, gains roughly one polynomial degree
        of definiteness (certified in N.7b), and combined with the chamfer
        (which caps $s$ at $\approx 0.7$ and removes two-face cells) is
        the predicted robust \emph{full-order} $p=4$ configuration.
\end{enumerate}

\subsection*{N.1 Setting}

Model: cell $K = (0,h)^d$, $\mathbb Q_p$ Lagrange polynomials, SBM form
restricted to $K$ and its surrogate faces $F$:
\begin{equation}
  a_K(u,v) = \int_K \nabla u\cdot\nabla v
  - \int_F (\partial_n u)\, v
  - \int_F (\partial_n v)\, \mathcal E u
  + \sigma \int_F \mathcal E u\, \mathcal E v,
  \qquad \sigma = \gamma\, p^2/h,
  \label{eq:cellform}
\end{equation}
with $\mathcal E u(\mathbf x) = \sum_{j=0}^{k} \tfrac1{j!} (\mathbf
d\cdot\nabla)^j u(\mathbf x)$. Throughout, $s = \|\mathbf d\|/h$ and we set
$h=1$ in computations; $\gamma = 4$ unless stated (results are insensitive
to $\gamma$, see N.4).

Two mechanisms could plausibly cause the observed breakdown: (i) elliptic
corner singularities of the surrogate domain (reentrant staircase corners,
interior angle $3\pi/2$, Kondrat'ev exponent $r^{2/3}$), and (ii) loss of
definiteness of the discrete boundary terms. The results below show that
(ii) is dominant and quantitatively predictive; (i) is discussed briefly in
N.8.

\subsection*{N.2 Sharp extrapolation amplification}

\paragraph{Lemma N.1 (extension = extrapolation; sharp growth).}
\emph{Let $q$ be a polynomial of degree $p$ on $[0,h]$ and let the shift be
axis-aligned. For $k \ge p$ the Taylor extension is exact point evaluation,
$\mathcal E q(0) = q(-\delta)$, and}
\[
  \sup\Big\{ |q(-\delta)| \,:\, \|q\|_{L^\infty(0,h)} \le 1 \Big\}
  \;=\; T_p\!\big(1 + 2s\big), \qquad s = \delta/h,
\]
\emph{attained by the Chebyshev polynomial. Moreover, with
$\rho(s) := \sqrt s + \sqrt{1+s}$,}
\[
  T_p(1+2s) \;=\; \tfrac12\Big( \rho(s)^{2p} + \rho(s)^{-2p} \Big).
\]

\emph{Proof.} Exactness: the Taylor polynomial of order $k\ge p$ of a
degree-$p$ polynomial is the polynomial itself. Extremality of $T_p$
outside $[-1,1]$ is classical (Chebyshev; see Rivlin). The identity follows
from $T_p(\cosh\theta) = \cosh(p\theta)$ with
$1+2s = \cosh\big(2\,\mathrm{arcsinh}\sqrt s\big)$ and
$e^{\mathrm{arcsinh}\sqrt s} = \sqrt s + \sqrt{1+s}$. $\square$

The quantity $\rho(s)^2$ is the \emph{amplification per polynomial degree}:

\begin{center}
\begin{tabular}{lccccc}
\hline
 & $s=\tfrac14$ & $s=\tfrac12$ & $s=\tfrac34$ & $s=1$ & $s=\sqrt2$ \\
\hline
$\rho(s)^2$ (per degree)      & $2.6$ & $3.7$ & $4.8$ & $5.8$ & $7.5$ \\
$T_1(1+2s)$ & $1.5$ & $2$ & $2.5$ & $3$ & $3.8$ \\
$T_2(1+2s)$ & $3.5$ & $7$ & $11.5$ & $17$ & $28$ \\
$T_3(1+2s)$ & $9$ & $26$ & $55$ & $99$ & $214$ \\
$T_4(1+2s)$ & $24.5$ & $97$ & $265$ & $577$ & $1611$ \\
\hline
\end{tabular}
\end{center}

The penalty and consistency terms in \eqref{eq:cellform} are bilinear in
$\mathcal E$, so the boundary block scales like $\sigma\, T_p^2$ relative to
the $O(1)$ (in $h^{d-2}$ units) volume stiffness. This is confirmed by the
computed largest eigenvalues $\lambda_{\max}$ of the symmetric part: e.g.\
for the flat configuration at $p=4$, the measured ratios
$\lambda_{\max}(s{=}1)/\lambda_{\max}(s{=}\tfrac12) = 34$ and
$\lambda_{\max}(s{=}\sqrt2)/\lambda_{\max}(s{=}1) = 7.6$ match the predicted
$\big(T_4(3)/T_4(2)\big)^2 = 35$ and $\big(T_4(3.83)/T_4(3)\big)^2 = 7.8$.
At $p=4$, $s=1$ the local block condition number already exceeds $10^7$.
Conditioning alone, however, is not the failure mode --- definiteness is,
as the next two subsections show.

\subsection*{N.3 Exact one-dimensional results}

On $K=(0,h)$ with surrogate face at $x=0$ and true boundary at $-\delta$
(outward shift), $\mathcal E v(0) = v(-\delta)$, form \eqref{eq:cellform}.

\paragraph{Lemma N.2 ($p=1$, closed form).} \emph{For $v = a + bx$:}
\[
  a_K(v,v) = b^2 (h+\delta) + 2bw + \sigma w^2, \qquad w := v(-\delta),
\]
\emph{which is positive definite iff $\sigma (h + \delta) > 1$. For the
inward shift ($\Gamma$ at $+\delta$ inside the cell) the same computation
gives $b^2(h-\delta) + 2bw + \sigma w^2$, positive definite iff
$\sigma (h - \delta) > 1$.}

Consequences: $p=1$ with outward shifts is unconditionally coercive (any
$\sigma \ge 1/h$, any $\delta > 0$) --- matching the total robustness of
$p=1$ in practice. Inward shifts lose coercivity as $\delta \to h$; with
$\sigma = 4/h$ the threshold is $\delta = 3h/4$, reproduced to machine
precision by the 2D code (a useful cross-validation of both).

\paragraph{Lemma N.3 (kernel escape).} \emph{On the subspace
$\{v : \mathcal E v = 0 \text{ on } F\}$ the penalty and one flux term
vanish identically and \eqref{eq:cellform} reduces to the
\emph{unpenalized half-Nitsche} form}
\[
  F(v) = \|\nabla v\|_{L^2(K)}^2 - \int_F (\partial_n v)\, v .
\]
Definiteness on this kernel is a sign-correlation property between $v$ and
$\partial_n v$ on $F$, forced by $v$ vanishing at the shifted points; the
penalty is powerless there, no matter how large. This is the structural
reason for the $\gamma$-insensitivity observed in N.4.

\paragraph{Lemma N.4 ($p=2$, 1D, explicit threshold).} \emph{For $v$
quadratic with $v(-\delta) = 0$ (kernel of $\mathcal E$), writing
$c = v'(0) $-parametrization $v = (x+\delta)(a+bx)$, $h=1$:}
\[
  F(v) = (1+s)\,c^2 + (2 - s^2)\,cb + \tfrac43\, b^2, \qquad c = a + bs,
\]
\emph{positive definite iff $(2-s^2)^2 < \tfrac{16}{3}(1+s)$, i.e.\ for all
$s \lesssim 2.5$.}

So in 1D even $p=2$ survives shifts of two and a half cells; the observed
2D breakdown at $s \le \sqrt2$ cannot be a 1D phenomenon. What 1D misses:
at a staircase step, \emph{two} faces with orthogonal shift directions load
the same cell polynomial, and the sign-correlation of Lemma N.3 cannot be
satisfied for both directions simultaneously once $p$ gives the polynomial
enough freedom to oscillate between the cell and the shifted evaluation
points.

\subsection*{N.4 Numerical definiteness certificates in 2D}

\paragraph{Certificate construction.} Let $K$ be a boundary cell in one of
the configurations below and let $v \in \mathbb Q_p(K)$ vanish on the edges
of $K$ shared with active neighbors. Extending $v$ by zero yields a
conforming global function whose global energy equals $a_K(v,v)$: volume
terms of neighbors vanish, and all other surrogate-face terms vanish
because $\mathcal E$ on a face of cell $K'$ uses only the polynomial of
$K'$. Hence
\[
  \lambda_{\min}\Big( \tfrac12 (A_K + A_K^T)\big|_{S} \Big) < 0
  \;\Longrightarrow\;
  \text{the \emph{global} symmetric part is indefinite}
\]
on every mesh containing the configuration, where $S$ is the zero-extension
subspace. Configurations ($h=1$, exact extrapolation $=$ Taylor with $k=p$
for the axis-aligned shifts):
\emph{flat} (one face $\{y{=}1\}$, $\mathbf d = (0,s)$; $S$: $v=0$ on the
three interior edges), \emph{step} (two faces $\{y{=}1\},\{x{=}1\}$,
$\mathbf d = (0,s),(s,0)$; the staircase step cell; $S$: $v=0$ on two
interior edges), \emph{step-diag} (as step but both shifts
$s(1,1)/\sqrt2$, modeling pocket-directed closest-point projections), and
\emph{chamfer} variants in N.6.

\paragraph{Results ($\gamma = 4$).} Minimal eigenvalue of the symmetric
part on $S$ (negative $=$ certified global indefiniteness; the volume
stiffness is $O(1)$ in these units, so magnitudes are meaningful):

\begin{center}
\begin{tabular}{llcccc}
\hline
config & $p$ & $s=0.5$ & $s=0.75$ & $s=1$ & $s=\sqrt2$ \\
\hline
flat  & 3 & $+1.22$ & $+0.91$ & $+0.75$ & $+0.61$ \\
flat  & 4 & $+0.28$ & $+0.08$ & $\mathbf{-0.03}$ & $\mathbf{-0.13}$ \\
step  & 3 & $+0.27$ & $+0.09$ & $+0.02$ & $\mathbf{-0.05}$ \\
step  & 4 & $\mathbf{-0.04}$ & $\mathbf{-0.13}$ & $\mathbf{-0.18}$ & $\mathbf{-0.24}$ \\
\hline
\end{tabular}
\end{center}

Findings:
\begin{enumerate}
  \item[(i)] $p \le 2$: positive in every configuration and shift tested
        (not shown; margins $O(1)$).
  \item[(ii)] $p=3$: indefinite \emph{only} at the two-face step cell with
        $s \to \sqrt2$, i.e.\ at deep reentrant pockets. Whether a given
        geometry produces such pockets depends on how $\Gamma$ cuts the
        grid --- hence ``$p=3$ sometimes breaks''.
  \item[(iii)] $p=4$: indefinite from $s \approx 0.5$ at step cells and
        from $s \approx 1$ at flat faces. Every staircase surrogate of a
        generic $\Gamma$ contains such cells --- hence ``$p=4$ always
        breaks''.
  \item[(iv)] The nonsymmetric spectra show only small imaginary parts for
        outward shifts: for $\lambda = 0$ the pathology is
        \emph{indefiniteness of the symmetric part}, not complex
        eigenvalues (those belong to inward shifts, consistent with the 1D
        study in~\cite{wichrowski2025geometric}).
  \item[(v)] Penalty insensitivity (Lemma N.3 in action), step cell,
        $s=1$: for $p=4$, $\lambda_{\min}|_S = -0.191, -0.184, -0.178,
        -0.177$ at $\gamma = 2, 4, 16, 64$. No penalty tuning can repair
        this.
  \item[(vi)] Inward shifts (flat face, $\mathbf d = (0,-s)$, corrected
        certificate): $p=2$ positive at all $s$; $p=3$ and $p=4$
        \emph{negative at every} $s \in \{0.25, 0.5, 0.75\}$ (e.g.\ $-0.63$
        and $-1.13$ at $s=0.25$). Intersected-cell surrogates
        ($\lambda>0$) are therefore incompatible with $p \ge 3$
        independently of corner effects.
\end{enumerate}

\subsection*{N.5 Consequences for smoothers}

Let $v_\star$ be a certified negative mode ($a(v_\star,v_\star) < 0$,
supported in one boundary cell).

\begin{itemize}
  \item \textbf{All energy-based convergence theory is void.} Both the
        Schwarz theory for additive/multiplicative smoothers and the
        field-of-values GMRES bounds (Eisenstat--Elman--Schultz,
        Cai--Widlund) presuppose a definite symmetric part. On staircase
        SBM with $p \ge 3$ this hypothesis \emph{fails as a matter of
        fact}, not of proof technique.
  \item \textbf{Pointwise and cell-wise additive smoothers.} For
        (block-)Jacobi-type methods the local blocks themselves become
        indefinite (the single-cell Neumann blocks go negative one degree
        earlier than the certificates, already $p=3$, $s \ge 0.75$).
        Damped Richardson with an indefinite block-diagonal preconditioner
        has no admissible damping: the modes with negative Rayleigh
        quotient are amplified for every $\omega > 0$ of either sign
        globally. This is the cleanest available explanation for the
        observed \emph{divergence} (not merely stagnation) of the
        cell-wise additive DG+SBM smoother.
  \item \textbf{Why multiplicative full-residual patches still work at
        $p \le 3$.} The negative mode $v_\star$ is supported in one cell,
        hence entirely inside at least one vertex patch (the full-residual
        construction guarantees its whole residual row is visible there).
        An \emph{exact} local solve needs only invertibility, not
        definiteness, of $A_i$; visiting that patch removes $v_\star$'s
        error component exactly, and the sequential sweep propagates the
        repaired residual outward. In the additive version the $2^d$
        overlapping patches containing $v_\star$ each compute a correction
        against an indefinite, non-normal block and their \emph{sum} is
        applied blindly; there is no self-correction within the sweep, and
        the usual multiplicity damping $\omega \le 1/N_c$ cannot produce
        contraction where the field of values crosses zero.
  \item \textbf{Why $p=4$ defeats even the multiplicative smoother.} The
        indefinite modes are no longer isolated: at $p=4$ they exist at
        \emph{every} boundary cell with $s \gtrsim 0.5$--$1$, they couple
        to neighboring patches through off-block entries of size
        $O(\sigma T_p^2)$ ($\ge 10^6$ at $s = 1$), and the error injected
        across a patch boundary is re-amplified faster than the sweep
        damps it. Locality of the repair mechanism is lost.
\end{itemize}

\subsection*{N.6 Chamfer (mesh-defined triangular cuts): rescues $p=3$, not $p=4$}

The chamfer splits staircase corner quads along the vertex-to-vertex
diagonal. The cut position is fixed by the mesh: every cut cell is exactly
half a quad, uniformly shape-regular --- \emph{there is no small-cut/sliver
problem and no ghost penalty is needed}, in contrast to
$\Gamma$-interpolating CutFEM cuts. Reentrant $3\pi/2$ corners become pairs
of $5\pi/4$ corners (regularity exponent $2/3 \to 4/5$), and the worst
shift drops from $s \approx \sqrt2$ to $s \approx 0.7$ with a single,
near-normal face replacing the two-face loading.

Certificates for the chamfered cell (hypotenuse face, shift
$s(1,1)/\sqrt2$), two variants: $\mathbb P_p$ on the triangle (remeshed)
and $\mathbb Q_p$ on the parent quad integrated over the triangle
(cut-quadrature):

\begin{center}
\begin{tabular}{llccc}
\hline
variant & $p$ & $s=0.5$ & $s=0.707$ & $s=1$ \\
\hline
$\mathbb P_p$ triangle & 3 & $+0.91$ & $+0.62$ & $+0.42$ \\
$\mathbb P_p$ triangle & 4 & $\mathbf{-0.15}$ & $\mathbf{-0.34}$ & $\mathbf{-0.48}$ \\
$\mathbb Q_p$ cut      & 3 & $\mathbf{-0.01}$ & $\approx\mathbf{-0.03}$ & $\mathbf{-0.02}$ \\
$\mathbb Q_p$ cut      & 4 & $\mathbf{-0.08}$ & $\mathbf{-0.07}$ & $\mathbf{-0.06}$ \\
\hline
\end{tabular}
\end{center}

Reading:
\begin{itemize}
  \item The \emph{remeshed} chamfer ($\mathbb P_p$ triangles) is positive
        for $p=3$ at all relevant shifts. Since the chamfer also caps $s$
        at $\approx 0.7$, it removes exactly the configuration in which
        $p=3$ was certified indefinite. \textbf{Prediction: chamfering
        makes $p=3$ uniformly robust.}
  \item It does \emph{not} rescue $p=4$: the triangle block is indefinite
        from $s = 0.5$. Two reasons: (a) a triangle is thin in the
        extrapolation direction (extent $h/\sqrt2$ normal to the
        hypotenuse, tapering), so the effective Chebyshev argument is
        larger than $s$ suggests; (b) independently, $p=4$ is already
        indefinite at \emph{flat} faces with $s \approx 1$, which the
        chamfer does not touch.
  \item The cut-quadrature variant is weaker (marginally indefinite even
        at $p=3$): the polynomial is energy-controlled only on the
        half-cell while being extrapolated from its edge. If chamfering is
        implemented, the triangles should carry genuine $\mathbb P_p$
        spaces (or the parent-cell variant needs a stabilization on the
        inactive half). This is an implementation-relevant distinction the
        earlier draft of these notes missed.
\end{itemize}

\subsection*{N.7 The effective fix: truncated extrapolation order}

Lemma N.1 identifies the amplification driver as the extrapolation
\emph{order}: truncating the Taylor extension at order $k < p$ caps the
growth at degree-$k$ extrapolation ($T_k$-like) while retaining
$O((sh)^{k+1})$ boundary consistency. Certificates for the step cell
(worst outward configuration), truncation order $k$:

\begin{center}
\begin{tabular}{llcccc}
\hline
config & $p$ & $k=p$ & $k=3$ & $k=2$ & $k=1$ \\
\hline
step, $s=0.75$   & 4 & $-0.13$ & $-0.10$ & $\mathbf{+0.01}$ & $+0.41$ \\
step, $s=1$      & 4 & $-0.18$ & $-0.14$ & $-0.03$ & $+0.35$ \\
step, $s=\sqrt2$ & 4 & $-0.24$ & $-0.18$ & $-0.06$ & $+0.30$ \\
step, $s=\sqrt2$ & 3 & $-0.05$ & --      & $\mathbf{+0.04}$ & $+0.50$ \\
flat, $s=1$      & 4 & $-0.03$ & $+0.06$ & $+0.34$ & $+0.98$ \\
\hline
\end{tabular}
\end{center}

\begin{itemize}
  \item $k=1$ restores definiteness \emph{everywhere} with $O(1)$ margins
        --- a bulletproof but low-order option (locally $O(h^2)$
        consistency; acceptable if applied only at the few worst faces).
  \item $k=2$ restores definiteness for $p=4$ up to $s \approx 0.75$ and
        for $p=3$ everywhere, with $O(h^3)$ boundary consistency.
  \item Honest objection: with a staircase or chamfered surrogate,
        $\delta = O(h)$, so $k<p$ caps the \emph{global} boundary accuracy
        at $O(h^{k+1})$ and largely defeats the purpose of $p=4$. (The
        truncation is order-neutral only where $\delta = O(h^2)$, i.e.\
        with a $\Gamma$-interpolating cut.) A shift-adaptive rule
        ($k=p$ where $s \le 0.4$, smaller $k$ beyond) confines the loss to
        the worst faces, but the principled resolution is to change the
        extension, not its order --- see N.7b.
  \item For inward shifts truncation helps at moderate $s$ but $p \ge 3$
        remains indefinite at small inward $s$ regardless of $k$ ---
        reinforcing finding (vi): use $\lambda = 0$ for $p \ge 3$.
\end{itemize}

\subsection*{N.7b Alternative extension operators: a no-go lemma and three ways around it}

\paragraph{Lemma N.5 (no free lunch for cell-local extensions).}
\emph{Let $\mathcal L : \mathbb Q_p(K) \to \mathbb R$ be linear, computed
from data on $K$ only, and exact on polynomials,
$\mathcal L q = q(\tilde{\mathbf x} + \mathbf d)$ for all $q \in
\mathbb Q_p(K)$ (as required for $O(h^{p+1})$ consistency from cell-local
data). Then $\mathcal L$ \emph{is} the point-evaluation extrapolation on
$\mathbb Q_p(K)$, and its norm is the Chebyshev factor of Lemma N.1.}
\emph{Proof:} two linear functionals agreeing on the whole space are
equal. $\square$

Smoothness of the extension is a red herring: on the discrete space every
exact extension coincides with the one already in use. Any genuine
alternative must therefore use information beyond the single cell's
polynomial. Exactly three sources exist, and each gives a known method
family:

\paragraph{(A) Use the PDE: substituted normal derivatives (recommended).}
Since $-\Delta u = f$, all normal derivatives of order $\ge 2$ can be
eliminated recursively,
\[
  \partial_n^2 u = -f - \Delta_{\mathrm{tan}} u, \qquad
  \partial_n^3 u = -\partial_n f - \Delta_{\mathrm{tan}} \partial_n u,
  \quad \dots
\]
so the degree-$k$ Taylor extension becomes
\[
  \mathcal E_{\mathrm{pde}} u
  = u + \delta\, \partial_n u
  + \sum_{j=2}^{k} \tfrac{\delta^j}{j!}
    \big( \text{tangential derivatives of } u,\ \partial_n u \big)
  + \underbrace{\big( f\text{-terms} \big)}_{\text{data, rhs}} ,
\]
with \emph{every $u$-dependent term evaluated on the face}: no
extrapolation beyond the first normal derivative. Crucially,
$\mathcal E_{\mathrm{pde}}$ is exact on solutions of the PDE --- which is
all that consistency requires --- so \textbf{the full order $O(\delta^{k+1})$
is retained} while the Chebyshev growth is traded for tangential Markov
factors $\big(c\,p^2 \delta/h\big)^{j}$, small for moderate $s$.
Certificates (same setup as N.4, $f$-part irrelevant for the spectrum):

\begin{center}
\begin{tabular}{llcccc}
\hline
config & $p$ & $s=0.5$ & $s=0.75$ & $s=1$ & $s=\sqrt2$ \\
\hline
flat, $\mathcal E_{\mathrm{pde}}$ & 3 & $+2.87$ & $+3.36$ & $+3.84$ & $+4.25$ \\
flat, $\mathcal E_{\mathrm{pde}}$ & 4 & $+1.34$ & $+0.80$ & $-0.01$ & $-0.71$ \\
step, $\mathcal E_{\mathrm{pde}}$ & 3 & $+1.20$ & $+0.90$ & $+0.49$ & $-0.10$ \\
step, $\mathcal E_{\mathrm{pde}}$ & 4 & $+0.12$ & $-0.34$ & $-0.38$ & $-0.33$ \\
\hline
\end{tabular}
\end{center}

Compared with plain Taylor (N.4): roughly one polynomial degree of
definiteness is gained at full consistency order ($p=4$ flat: $+0.08 \to
+0.80$ at $s = 0.75$; step: indefinite from $s=0.5 \to$ definite at
$s=0.5$). At $s \gtrsim 1$ the tangential factors take over and
indefiniteness returns --- so $\mathcal E_{\mathrm{pde}}$ alone does not
cure the staircase, but \textbf{chamfer ($s \le 0.7$, single near-normal
face) $+$ $\mathcal E_{\mathrm{pde}}$ is certified-definite in every
configuration that then occurs, at full order.} Implementation cost:
tangential/mixed derivatives of the shape functions on the face (already
available), plus $f$ and its normal-derivative data near the boundary in
the right-hand side.

\paragraph{(B) Use the boundary data: interpolatory (ghost/ROD-type)
conditions.} Taylor never uses the one thing that is known exactly: $g$ at
the projected point $\mathbf x \in \Gamma$. Writing the boundary condition
via a 1D interpolant along the shift line through $p{+}1$ points
\emph{inside} $\tilde\Omega$ plus the pinned value $g(\mathbf x)$, the
surrogate-face value $u(\tilde{\mathbf x})$ is expressed by
\emph{interpolation} (the evaluation point lies between the samples and
$\mathbf x$), so the $u$-dependent weights are bounded by a Lebesgue
constant instead of $T_p$; full order is retained. This is the
finite-element analogue of ghost-point methods (Gibou--Fedkiw) and
reconstructed off-site data (ROD). Cost: the stencil crosses 1--2 interior
cells, so the sparsity pattern widens; the samples stay inside a vertex
patch, so the full-residual construction survives.

\paragraph{(C) Use a larger support: aggregated extension.} Extrapolate
from a least-squares polynomial fit over an interior patch of diameter
$L = 2h$--$3h$ (AgFEM-flavored). The amplification drops to
$T_p(1 + 2\delta/L)$; by the certificate tables, an effective shift
$\delta/L \le 0.4$ is safe even at $p=4$, so $L \ge 3h$ suffices for the
worst staircase pockets. Cost: same widened coupling as (B), plus the fit;
conceptually the bluntest of the three.

Ranking for this project: (A) is the least invasive (face-local, no
sparsity change, no new data structures) and combines multiplicatively
with the chamfer; (B) is the most robust in principle (bounded weights
regardless of geometry) and is the natural fallback if (A)$+$chamfer still
misbehaves at $p \ge 5$ or in 3D, where edge/vertex pockets give $s$ up to
$\sqrt3$.

\subsection*{N.8 The regularity mechanism (demoted to second order)}

Reentrant staircase corners carry Kondrat'ev singularities
$r^{2/3}\sin(2\theta/3)$; they are recreated at new positions on every
level (non-nested active sets), so the associated error is $O(h)$-wavelength
by construction and falls entirely on the fine-level smoother, and
degree-$p$ approximation saturates at $O(h^{2/3})$ locally. This degrades
constants and localizes the hard error components at the boundary --- it
explains \emph{where} the smoother must work hardest, but not the observed
$p$-thresholds: the certificates of N.4 do, quantitatively and with the
right $\lambda$-, $s$-, and configuration-dependence. If the exact SBM
extrapolation were used ($k$ exact, no indefiniteness), the corner
singularities alone would be handled by the existing patch smoother, as
they are for $p \le 2$.

\subsection*{N.9 What is proven, what is certified, what is conjectured}

\emph{Proven} (Lemmas N.1--N.4): sharp Chebyshev amplification with the
per-degree factor $\rho(s)^2$; unconditional 1D $p=1$ coercivity (outward)
and the inward threshold $\sigma(h-\delta) > 1$; the kernel-escape
mechanism and penalty-insensitivity; 1D $p=2$ coercivity to $s\approx2.5$.

\emph{Certified numerically} (exact statements about finite matrices,
N.4/N.6/N.7): global indefiniteness of the symmetric part for the listed
$(p, s, \text{config})$; its penalty-insensitivity; positivity of the
chamfered $\mathbb P_3$ and truncated-$k$ variants.

\emph{Proven} (Lemma N.5): no cell-local exact extension can beat the
Chebyshev amplification; alternatives must use the PDE, the data $g$, or a
larger support (N.7b).

\emph{Conjectured / to validate in the solver}: (a) chamfer $\Rightarrow$
$p=3$ uniformly robust; (b) chamfer $+$ PDE-substituted extension
$\mathcal E_{\mathrm{pde}}$ $\Rightarrow$ $p=4$ robust \emph{at full
order} (the leading recommendation); (c) shift-adaptive $k$-truncation
alone (no chamfer) already recovers $p=4$ convergence with the existing
smoother at reduced boundary accuracy --- still the cheapest diagnostic
experiment: a ten-line change that directly tests the whole theory;
(d) restoring definiteness restores the additive patch smoother to
usefulness (with standard multiplicity damping), enabling the parallel
variant.

% --- references to add if these notes graduate into the paper:
% Rivlin, "Chebyshev Polynomials" (extremal growth outside [-1,1]).
% Grisvard, "Elliptic Problems in Nonsmooth Domains" (corner asymptotics).
% Atallah, Canuto, Scovazzi (SBM analysis, d/h conditions).
% Cai & Widlund SINUM 1992/93; Eisenstat, Elman, Schultz 1983 (nonsymmetric bounds).
% Burman et al. (CutFEM/ghost penalty, for the contrast drawn in N.6).
 at the end of sbm_patch_mg.tex to compile.
% All numerical values reproducible via theory_certificates.py (numpy only).
% ==========================================================================

\section*{Appendix (working notes): definiteness of the SBM boundary terms at the surrogate staircase, and consequences for smoothers}

\subsection*{N.0 Executive summary}

Write $s := \|\mathbf d\|/h$ for the shift length in units of the mesh size.
On a staircase surrogate ($\lambda = 0$) the flat faces have $s \lesssim 1$
and the cells at staircase steps carry \emph{two} extrapolated faces with
$s$ up to $\sqrt 2$ at reentrant pockets. The findings of these notes:

\begin{enumerate}
  \item The Taylor extension of a degree-$p$ polynomial over distance $sh$
        is polynomial extrapolation; its sharp amplification is the
        Chebyshev factor $T_p(1+2s)$, which grows by a factor
        $\rho(s)^2 = \big(\sqrt s + \sqrt{1+s}\big)^2$ \emph{per polynomial
        degree} ($\approx 3.7$ at $s=\tfrac12$, $\approx 7.5$ at
        $s=\sqrt2$). This scaling is confirmed to within a few percent by
        the computed spectra (N.4).
  \item In 1D the method is provably healthy: $p=1$ is unconditionally
        coercive for outward shifts (Lemma N.2), and $p=2$ remains coercive
        up to the absurd shift $s \approx 2.5$ (Lemma N.4). \emph{The
        breakdown is intrinsically multi-dimensional.}
  \item Computable \emph{global indefiniteness certificates} (N.4): the
        symmetric part of the SBM bilinear form is indefinite ---
        on functions supported in a single boundary cell, hence on any mesh
        containing that configuration --- for
        \[
          p = 3:\ \text{only at two-face step cells with } s \to \sqrt 2
          \quad\text{(deep reentrant pockets)},
        \]
        \[
          p = 4:\ \text{from } s \approx 0.5 \text{ at step cells, } s
          \approx 1 \text{ even at flat faces.}
        \]
        This reproduces the empirical pattern exactly: $p\le2$ robust,
        $p=3$ fails \emph{sometimes} (only when the geometry creates deep
        pockets), $p=4$ fails \emph{always} (every staircase contains
        $s \ge 0.5$ step cells). Increasing the penalty $\sigma$ does
        \textbf{not} help (N.4, item v): the offending modes nearly
        annihilate $\mathcal E$, so they do not feel the penalty.
  \item Inward shifts ($\lambda > 0$, true boundary inside the surrogate)
        are far worse: indefinite for all $p \ge 3$ at \emph{every} shift
        length tested, and not repaired by any of the fixes below. For
        $p \ge 3$ the threshold $\lambda$ should be $0$.
  \item Chamfering the staircase with triangles (the proposal of these
        notes' earlier draft) \textbf{rescues $p=3$ but not $p=4$}: it
        removes exactly the deep-pocket configurations where $p=3$ fails,
        but $p=4$ is already indefinite at flat faces with $s \approx 1$,
        which chamfering does not shorten. There is no sliver problem
        (cuts are mesh-defined, all cells uniformly shape-regular), but a
        triangle is thin in the extrapolation direction, which eats part
        of the benefit (N.6).
  \item Truncating the Taylor order to $k=2$ ($k=1$) restores definiteness
        at step cells up to $s \approx 0.75$ (everywhere), at the price of
        reduced boundary consistency $O(h^{k+1})$ --- which defeats the
        purpose of $p=4$ (N.7).
  \item A no-go lemma (N.7b): \emph{any} linear cell-local extension exact
        on $\mathbb Q_p(K)$ coincides with polynomial extrapolation and
        carries the full Chebyshev amplification. Alternatives must import
        extra information: the PDE, the boundary data $g$, or a larger
        support. The PDE-substituted extension (all normal derivatives of
        order $\ge 2$ traded for tangential derivatives and $f$-data)
        keeps full consistency order, gains roughly one polynomial degree
        of definiteness (certified in N.7b), and combined with the chamfer
        (which caps $s$ at $\approx 0.7$ and removes two-face cells) is
        the predicted robust \emph{full-order} $p=4$ configuration.
\end{enumerate}

\subsection*{N.1 Setting}

Model: cell $K = (0,h)^d$, $\mathbb Q_p$ Lagrange polynomials, SBM form
restricted to $K$ and its surrogate faces $F$:
\begin{equation}
  a_K(u,v) = \int_K \nabla u\cdot\nabla v
  - \int_F (\partial_n u)\, v
  - \int_F (\partial_n v)\, \mathcal E u
  + \sigma \int_F \mathcal E u\, \mathcal E v,
  \qquad \sigma = \gamma\, p^2/h,
  \label{eq:cellform}
\end{equation}
with $\mathcal E u(\mathbf x) = \sum_{j=0}^{k} \tfrac1{j!} (\mathbf
d\cdot\nabla)^j u(\mathbf x)$. Throughout, $s = \|\mathbf d\|/h$ and we set
$h=1$ in computations; $\gamma = 4$ unless stated (results are insensitive
to $\gamma$, see N.4).

Two mechanisms could plausibly cause the observed breakdown: (i) elliptic
corner singularities of the surrogate domain (reentrant staircase corners,
interior angle $3\pi/2$, Kondrat'ev exponent $r^{2/3}$), and (ii) loss of
definiteness of the discrete boundary terms. The results below show that
(ii) is dominant and quantitatively predictive; (i) is discussed briefly in
N.8.

\subsection*{N.2 Sharp extrapolation amplification}

\paragraph{Lemma N.1 (extension = extrapolation; sharp growth).}
\emph{Let $q$ be a polynomial of degree $p$ on $[0,h]$ and let the shift be
axis-aligned. For $k \ge p$ the Taylor extension is exact point evaluation,
$\mathcal E q(0) = q(-\delta)$, and}
\[
  \sup\Big\{ |q(-\delta)| \,:\, \|q\|_{L^\infty(0,h)} \le 1 \Big\}
  \;=\; T_p\!\big(1 + 2s\big), \qquad s = \delta/h,
\]
\emph{attained by the Chebyshev polynomial. Moreover, with
$\rho(s) := \sqrt s + \sqrt{1+s}$,}
\[
  T_p(1+2s) \;=\; \tfrac12\Big( \rho(s)^{2p} + \rho(s)^{-2p} \Big).
\]

\emph{Proof.} Exactness: the Taylor polynomial of order $k\ge p$ of a
degree-$p$ polynomial is the polynomial itself. Extremality of $T_p$
outside $[-1,1]$ is classical (Chebyshev; see Rivlin). The identity follows
from $T_p(\cosh\theta) = \cosh(p\theta)$ with
$1+2s = \cosh\big(2\,\mathrm{arcsinh}\sqrt s\big)$ and
$e^{\mathrm{arcsinh}\sqrt s} = \sqrt s + \sqrt{1+s}$. $\square$

The quantity $\rho(s)^2$ is the \emph{amplification per polynomial degree}:

\begin{center}
\begin{tabular}{lccccc}
\hline
 & $s=\tfrac14$ & $s=\tfrac12$ & $s=\tfrac34$ & $s=1$ & $s=\sqrt2$ \\
\hline
$\rho(s)^2$ (per degree)      & $2.6$ & $3.7$ & $4.8$ & $5.8$ & $7.5$ \\
$T_1(1+2s)$ & $1.5$ & $2$ & $2.5$ & $3$ & $3.8$ \\
$T_2(1+2s)$ & $3.5$ & $7$ & $11.5$ & $17$ & $28$ \\
$T_3(1+2s)$ & $9$ & $26$ & $55$ & $99$ & $214$ \\
$T_4(1+2s)$ & $24.5$ & $97$ & $265$ & $577$ & $1611$ \\
\hline
\end{tabular}
\end{center}

The penalty and consistency terms in \eqref{eq:cellform} are bilinear in
$\mathcal E$, so the boundary block scales like $\sigma\, T_p^2$ relative to
the $O(1)$ (in $h^{d-2}$ units) volume stiffness. This is confirmed by the
computed largest eigenvalues $\lambda_{\max}$ of the symmetric part: e.g.\
for the flat configuration at $p=4$, the measured ratios
$\lambda_{\max}(s{=}1)/\lambda_{\max}(s{=}\tfrac12) = 34$ and
$\lambda_{\max}(s{=}\sqrt2)/\lambda_{\max}(s{=}1) = 7.6$ match the predicted
$\big(T_4(3)/T_4(2)\big)^2 = 35$ and $\big(T_4(3.83)/T_4(3)\big)^2 = 7.8$.
At $p=4$, $s=1$ the local block condition number already exceeds $10^7$.
Conditioning alone, however, is not the failure mode --- definiteness is,
as the next two subsections show.

\subsection*{N.3 Exact one-dimensional results}

On $K=(0,h)$ with surrogate face at $x=0$ and true boundary at $-\delta$
(outward shift), $\mathcal E v(0) = v(-\delta)$, form \eqref{eq:cellform}.

\paragraph{Lemma N.2 ($p=1$, closed form).} \emph{For $v = a + bx$:}
\[
  a_K(v,v) = b^2 (h+\delta) + 2bw + \sigma w^2, \qquad w := v(-\delta),
\]
\emph{which is positive definite iff $\sigma (h + \delta) > 1$. For the
inward shift ($\Gamma$ at $+\delta$ inside the cell) the same computation
gives $b^2(h-\delta) + 2bw + \sigma w^2$, positive definite iff
$\sigma (h - \delta) > 1$.}

Consequences: $p=1$ with outward shifts is unconditionally coercive (any
$\sigma \ge 1/h$, any $\delta > 0$) --- matching the total robustness of
$p=1$ in practice. Inward shifts lose coercivity as $\delta \to h$; with
$\sigma = 4/h$ the threshold is $\delta = 3h/4$, reproduced to machine
precision by the 2D code (a useful cross-validation of both).

\paragraph{Lemma N.3 (kernel escape).} \emph{On the subspace
$\{v : \mathcal E v = 0 \text{ on } F\}$ the penalty and one flux term
vanish identically and \eqref{eq:cellform} reduces to the
\emph{unpenalized half-Nitsche} form}
\[
  F(v) = \|\nabla v\|_{L^2(K)}^2 - \int_F (\partial_n v)\, v .
\]
Definiteness on this kernel is a sign-correlation property between $v$ and
$\partial_n v$ on $F$, forced by $v$ vanishing at the shifted points; the
penalty is powerless there, no matter how large. This is the structural
reason for the $\gamma$-insensitivity observed in N.4.

\paragraph{Lemma N.4 ($p=2$, 1D, explicit threshold).} \emph{For $v$
quadratic with $v(-\delta) = 0$ (kernel of $\mathcal E$), writing
$c = v'(0) $-parametrization $v = (x+\delta)(a+bx)$, $h=1$:}
\[
  F(v) = (1+s)\,c^2 + (2 - s^2)\,cb + \tfrac43\, b^2, \qquad c = a + bs,
\]
\emph{positive definite iff $(2-s^2)^2 < \tfrac{16}{3}(1+s)$, i.e.\ for all
$s \lesssim 2.5$.}

So in 1D even $p=2$ survives shifts of two and a half cells; the observed
2D breakdown at $s \le \sqrt2$ cannot be a 1D phenomenon. What 1D misses:
at a staircase step, \emph{two} faces with orthogonal shift directions load
the same cell polynomial, and the sign-correlation of Lemma N.3 cannot be
satisfied for both directions simultaneously once $p$ gives the polynomial
enough freedom to oscillate between the cell and the shifted evaluation
points.

\subsection*{N.4 Numerical definiteness certificates in 2D}

\paragraph{Certificate construction.} Let $K$ be a boundary cell in one of
the configurations below and let $v \in \mathbb Q_p(K)$ vanish on the edges
of $K$ shared with active neighbors. Extending $v$ by zero yields a
conforming global function whose global energy equals $a_K(v,v)$: volume
terms of neighbors vanish, and all other surrogate-face terms vanish
because $\mathcal E$ on a face of cell $K'$ uses only the polynomial of
$K'$. Hence
\[
  \lambda_{\min}\Big( \tfrac12 (A_K + A_K^T)\big|_{S} \Big) < 0
  \;\Longrightarrow\;
  \text{the \emph{global} symmetric part is indefinite}
\]
on every mesh containing the configuration, where $S$ is the zero-extension
subspace. Configurations ($h=1$, exact extrapolation $=$ Taylor with $k=p$
for the axis-aligned shifts):
\emph{flat} (one face $\{y{=}1\}$, $\mathbf d = (0,s)$; $S$: $v=0$ on the
three interior edges), \emph{step} (two faces $\{y{=}1\},\{x{=}1\}$,
$\mathbf d = (0,s),(s,0)$; the staircase step cell; $S$: $v=0$ on two
interior edges), \emph{step-diag} (as step but both shifts
$s(1,1)/\sqrt2$, modeling pocket-directed closest-point projections), and
\emph{chamfer} variants in N.6.

\paragraph{Results ($\gamma = 4$).} Minimal eigenvalue of the symmetric
part on $S$ (negative $=$ certified global indefiniteness; the volume
stiffness is $O(1)$ in these units, so magnitudes are meaningful):

\begin{center}
\begin{tabular}{llcccc}
\hline
config & $p$ & $s=0.5$ & $s=0.75$ & $s=1$ & $s=\sqrt2$ \\
\hline
flat  & 3 & $+1.22$ & $+0.91$ & $+0.75$ & $+0.61$ \\
flat  & 4 & $+0.28$ & $+0.08$ & $\mathbf{-0.03}$ & $\mathbf{-0.13}$ \\
step  & 3 & $+0.27$ & $+0.09$ & $+0.02$ & $\mathbf{-0.05}$ \\
step  & 4 & $\mathbf{-0.04}$ & $\mathbf{-0.13}$ & $\mathbf{-0.18}$ & $\mathbf{-0.24}$ \\
\hline
\end{tabular}
\end{center}

Findings:
\begin{enumerate}
  \item[(i)] $p \le 2$: positive in every configuration and shift tested
        (not shown; margins $O(1)$).
  \item[(ii)] $p=3$: indefinite \emph{only} at the two-face step cell with
        $s \to \sqrt2$, i.e.\ at deep reentrant pockets. Whether a given
        geometry produces such pockets depends on how $\Gamma$ cuts the
        grid --- hence ``$p=3$ sometimes breaks''.
  \item[(iii)] $p=4$: indefinite from $s \approx 0.5$ at step cells and
        from $s \approx 1$ at flat faces. Every staircase surrogate of a
        generic $\Gamma$ contains such cells --- hence ``$p=4$ always
        breaks''.
  \item[(iv)] The nonsymmetric spectra show only small imaginary parts for
        outward shifts: for $\lambda = 0$ the pathology is
        \emph{indefiniteness of the symmetric part}, not complex
        eigenvalues (those belong to inward shifts, consistent with the 1D
        study in~\cite{wichrowski2025geometric}).
  \item[(v)] Penalty insensitivity (Lemma N.3 in action), step cell,
        $s=1$: for $p=4$, $\lambda_{\min}|_S = -0.191, -0.184, -0.178,
        -0.177$ at $\gamma = 2, 4, 16, 64$. No penalty tuning can repair
        this.
  \item[(vi)] Inward shifts (flat face, $\mathbf d = (0,-s)$, corrected
        certificate): $p=2$ positive at all $s$; $p=3$ and $p=4$
        \emph{negative at every} $s \in \{0.25, 0.5, 0.75\}$ (e.g.\ $-0.63$
        and $-1.13$ at $s=0.25$). Intersected-cell surrogates
        ($\lambda>0$) are therefore incompatible with $p \ge 3$
        independently of corner effects.
\end{enumerate}

\subsection*{N.5 Consequences for smoothers}

Let $v_\star$ be a certified negative mode ($a(v_\star,v_\star) < 0$,
supported in one boundary cell).

\begin{itemize}
  \item \textbf{All energy-based convergence theory is void.} Both the
        Schwarz theory for additive/multiplicative smoothers and the
        field-of-values GMRES bounds (Eisenstat--Elman--Schultz,
        Cai--Widlund) presuppose a definite symmetric part. On staircase
        SBM with $p \ge 3$ this hypothesis \emph{fails as a matter of
        fact}, not of proof technique.
  \item \textbf{Pointwise and cell-wise additive smoothers.} For
        (block-)Jacobi-type methods the local blocks themselves become
        indefinite (the single-cell Neumann blocks go negative one degree
        earlier than the certificates, already $p=3$, $s \ge 0.75$).
        Damped Richardson with an indefinite block-diagonal preconditioner
        has no admissible damping: the modes with negative Rayleigh
        quotient are amplified for every $\omega > 0$ of either sign
        globally. This is the cleanest available explanation for the
        observed \emph{divergence} (not merely stagnation) of the
        cell-wise additive DG+SBM smoother.
  \item \textbf{Why multiplicative full-residual patches still work at
        $p \le 3$.} The negative mode $v_\star$ is supported in one cell,
        hence entirely inside at least one vertex patch (the full-residual
        construction guarantees its whole residual row is visible there).
        An \emph{exact} local solve needs only invertibility, not
        definiteness, of $A_i$; visiting that patch removes $v_\star$'s
        error component exactly, and the sequential sweep propagates the
        repaired residual outward. In the additive version the $2^d$
        overlapping patches containing $v_\star$ each compute a correction
        against an indefinite, non-normal block and their \emph{sum} is
        applied blindly; there is no self-correction within the sweep, and
        the usual multiplicity damping $\omega \le 1/N_c$ cannot produce
        contraction where the field of values crosses zero.
  \item \textbf{Why $p=4$ defeats even the multiplicative smoother.} The
        indefinite modes are no longer isolated: at $p=4$ they exist at
        \emph{every} boundary cell with $s \gtrsim 0.5$--$1$, they couple
        to neighboring patches through off-block entries of size
        $O(\sigma T_p^2)$ ($\ge 10^6$ at $s = 1$), and the error injected
        across a patch boundary is re-amplified faster than the sweep
        damps it. Locality of the repair mechanism is lost.
\end{itemize}

\subsection*{N.6 Chamfer (mesh-defined triangular cuts): rescues $p=3$, not $p=4$}

The chamfer splits staircase corner quads along the vertex-to-vertex
diagonal. The cut position is fixed by the mesh: every cut cell is exactly
half a quad, uniformly shape-regular --- \emph{there is no small-cut/sliver
problem and no ghost penalty is needed}, in contrast to
$\Gamma$-interpolating CutFEM cuts. Reentrant $3\pi/2$ corners become pairs
of $5\pi/4$ corners (regularity exponent $2/3 \to 4/5$), and the worst
shift drops from $s \approx \sqrt2$ to $s \approx 0.7$ with a single,
near-normal face replacing the two-face loading.

Certificates for the chamfered cell (hypotenuse face, shift
$s(1,1)/\sqrt2$), two variants: $\mathbb P_p$ on the triangle (remeshed)
and $\mathbb Q_p$ on the parent quad integrated over the triangle
(cut-quadrature):

\begin{center}
\begin{tabular}{llccc}
\hline
variant & $p$ & $s=0.5$ & $s=0.707$ & $s=1$ \\
\hline
$\mathbb P_p$ triangle & 3 & $+0.91$ & $+0.62$ & $+0.42$ \\
$\mathbb P_p$ triangle & 4 & $\mathbf{-0.15}$ & $\mathbf{-0.34}$ & $\mathbf{-0.48}$ \\
$\mathbb Q_p$ cut      & 3 & $\mathbf{-0.01}$ & $\approx\mathbf{-0.03}$ & $\mathbf{-0.02}$ \\
$\mathbb Q_p$ cut      & 4 & $\mathbf{-0.08}$ & $\mathbf{-0.07}$ & $\mathbf{-0.06}$ \\
\hline
\end{tabular}
\end{center}

Reading:
\begin{itemize}
  \item The \emph{remeshed} chamfer ($\mathbb P_p$ triangles) is positive
        for $p=3$ at all relevant shifts. Since the chamfer also caps $s$
        at $\approx 0.7$, it removes exactly the configuration in which
        $p=3$ was certified indefinite. \textbf{Prediction: chamfering
        makes $p=3$ uniformly robust.}
  \item It does \emph{not} rescue $p=4$: the triangle block is indefinite
        from $s = 0.5$. Two reasons: (a) a triangle is thin in the
        extrapolation direction (extent $h/\sqrt2$ normal to the
        hypotenuse, tapering), so the effective Chebyshev argument is
        larger than $s$ suggests; (b) independently, $p=4$ is already
        indefinite at \emph{flat} faces with $s \approx 1$, which the
        chamfer does not touch.
  \item The cut-quadrature variant is weaker (marginally indefinite even
        at $p=3$): the polynomial is energy-controlled only on the
        half-cell while being extrapolated from its edge. If chamfering is
        implemented, the triangles should carry genuine $\mathbb P_p$
        spaces (or the parent-cell variant needs a stabilization on the
        inactive half). This is an implementation-relevant distinction the
        earlier draft of these notes missed.
\end{itemize}

\subsection*{N.7 The effective fix: truncated extrapolation order}

Lemma N.1 identifies the amplification driver as the extrapolation
\emph{order}: truncating the Taylor extension at order $k < p$ caps the
growth at degree-$k$ extrapolation ($T_k$-like) while retaining
$O((sh)^{k+1})$ boundary consistency. Certificates for the step cell
(worst outward configuration), truncation order $k$:

\begin{center}
\begin{tabular}{llcccc}
\hline
config & $p$ & $k=p$ & $k=3$ & $k=2$ & $k=1$ \\
\hline
step, $s=0.75$   & 4 & $-0.13$ & $-0.10$ & $\mathbf{+0.01}$ & $+0.41$ \\
step, $s=1$      & 4 & $-0.18$ & $-0.14$ & $-0.03$ & $+0.35$ \\
step, $s=\sqrt2$ & 4 & $-0.24$ & $-0.18$ & $-0.06$ & $+0.30$ \\
step, $s=\sqrt2$ & 3 & $-0.05$ & --      & $\mathbf{+0.04}$ & $+0.50$ \\
flat, $s=1$      & 4 & $-0.03$ & $+0.06$ & $+0.34$ & $+0.98$ \\
\hline
\end{tabular}
\end{center}

\begin{itemize}
  \item $k=1$ restores definiteness \emph{everywhere} with $O(1)$ margins
        --- a bulletproof but low-order option (locally $O(h^2)$
        consistency; acceptable if applied only at the few worst faces).
  \item $k=2$ restores definiteness for $p=4$ up to $s \approx 0.75$ and
        for $p=3$ everywhere, with $O(h^3)$ boundary consistency.
  \item Honest objection: with a staircase or chamfered surrogate,
        $\delta = O(h)$, so $k<p$ caps the \emph{global} boundary accuracy
        at $O(h^{k+1})$ and largely defeats the purpose of $p=4$. (The
        truncation is order-neutral only where $\delta = O(h^2)$, i.e.\
        with a $\Gamma$-interpolating cut.) A shift-adaptive rule
        ($k=p$ where $s \le 0.4$, smaller $k$ beyond) confines the loss to
        the worst faces, but the principled resolution is to change the
        extension, not its order --- see N.7b.
  \item For inward shifts truncation helps at moderate $s$ but $p \ge 3$
        remains indefinite at small inward $s$ regardless of $k$ ---
        reinforcing finding (vi): use $\lambda = 0$ for $p \ge 3$.
\end{itemize}

\subsection*{N.7b Alternative extension operators: a no-go lemma and three ways around it}

\paragraph{Lemma N.5 (no free lunch for cell-local extensions).}
\emph{Let $\mathcal L : \mathbb Q_p(K) \to \mathbb R$ be linear, computed
from data on $K$ only, and exact on polynomials,
$\mathcal L q = q(\tilde{\mathbf x} + \mathbf d)$ for all $q \in
\mathbb Q_p(K)$ (as required for $O(h^{p+1})$ consistency from cell-local
data). Then $\mathcal L$ \emph{is} the point-evaluation extrapolation on
$\mathbb Q_p(K)$, and its norm is the Chebyshev factor of Lemma N.1.}
\emph{Proof:} two linear functionals agreeing on the whole space are
equal. $\square$

Smoothness of the extension is a red herring: on the discrete space every
exact extension coincides with the one already in use. Any genuine
alternative must therefore use information beyond the single cell's
polynomial. Exactly three sources exist, and each gives a known method
family:

\paragraph{(A) Use the PDE: substituted normal derivatives (recommended).}
Since $-\Delta u = f$, all normal derivatives of order $\ge 2$ can be
eliminated recursively,
\[
  \partial_n^2 u = -f - \Delta_{\mathrm{tan}} u, \qquad
  \partial_n^3 u = -\partial_n f - \Delta_{\mathrm{tan}} \partial_n u,
  \quad \dots
\]
so the degree-$k$ Taylor extension becomes
\[
  \mathcal E_{\mathrm{pde}} u
  = u + \delta\, \partial_n u
  + \sum_{j=2}^{k} \tfrac{\delta^j}{j!}
    \big( \text{tangential derivatives of } u,\ \partial_n u \big)
  + \underbrace{\big( f\text{-terms} \big)}_{\text{data, rhs}} ,
\]
with \emph{every $u$-dependent term evaluated on the face}: no
extrapolation beyond the first normal derivative. Crucially,
$\mathcal E_{\mathrm{pde}}$ is exact on solutions of the PDE --- which is
all that consistency requires --- so \textbf{the full order $O(\delta^{k+1})$
is retained} while the Chebyshev growth is traded for tangential Markov
factors $\big(c\,p^2 \delta/h\big)^{j}$, small for moderate $s$.
Certificates (same setup as N.4, $f$-part irrelevant for the spectrum):

\begin{center}
\begin{tabular}{llcccc}
\hline
config & $p$ & $s=0.5$ & $s=0.75$ & $s=1$ & $s=\sqrt2$ \\
\hline
flat, $\mathcal E_{\mathrm{pde}}$ & 3 & $+2.87$ & $+3.36$ & $+3.84$ & $+4.25$ \\
flat, $\mathcal E_{\mathrm{pde}}$ & 4 & $+1.34$ & $+0.80$ & $-0.01$ & $-0.71$ \\
step, $\mathcal E_{\mathrm{pde}}$ & 3 & $+1.20$ & $+0.90$ & $+0.49$ & $-0.10$ \\
step, $\mathcal E_{\mathrm{pde}}$ & 4 & $+0.12$ & $-0.34$ & $-0.38$ & $-0.33$ \\
\hline
\end{tabular}
\end{center}

Compared with plain Taylor (N.4): roughly one polynomial degree of
definiteness is gained at full consistency order ($p=4$ flat: $+0.08 \to
+0.80$ at $s = 0.75$; step: indefinite from $s=0.5 \to$ definite at
$s=0.5$). At $s \gtrsim 1$ the tangential factors take over and
indefiniteness returns --- so $\mathcal E_{\mathrm{pde}}$ alone does not
cure the staircase, but \textbf{chamfer ($s \le 0.7$, single near-normal
face) $+$ $\mathcal E_{\mathrm{pde}}$ is certified-definite in every
configuration that then occurs, at full order.} Implementation cost:
tangential/mixed derivatives of the shape functions on the face (already
available), plus $f$ and its normal-derivative data near the boundary in
the right-hand side.

\paragraph{(B) Use the boundary data: interpolatory (ghost/ROD-type)
conditions.} Taylor never uses the one thing that is known exactly: $g$ at
the projected point $\mathbf x \in \Gamma$. Writing the boundary condition
via a 1D interpolant along the shift line through $p{+}1$ points
\emph{inside} $\tilde\Omega$ plus the pinned value $g(\mathbf x)$, the
surrogate-face value $u(\tilde{\mathbf x})$ is expressed by
\emph{interpolation} (the evaluation point lies between the samples and
$\mathbf x$), so the $u$-dependent weights are bounded by a Lebesgue
constant instead of $T_p$; full order is retained. This is the
finite-element analogue of ghost-point methods (Gibou--Fedkiw) and
reconstructed off-site data (ROD). Cost: the stencil crosses 1--2 interior
cells, so the sparsity pattern widens; the samples stay inside a vertex
patch, so the full-residual construction survives.

\paragraph{(C) Use a larger support: aggregated extension.} Extrapolate
from a least-squares polynomial fit over an interior patch of diameter
$L = 2h$--$3h$ (AgFEM-flavored). The amplification drops to
$T_p(1 + 2\delta/L)$; by the certificate tables, an effective shift
$\delta/L \le 0.4$ is safe even at $p=4$, so $L \ge 3h$ suffices for the
worst staircase pockets. Cost: same widened coupling as (B), plus the fit;
conceptually the bluntest of the three.

Ranking for this project: (A) is the least invasive (face-local, no
sparsity change, no new data structures) and combines multiplicatively
with the chamfer; (B) is the most robust in principle (bounded weights
regardless of geometry) and is the natural fallback if (A)$+$chamfer still
misbehaves at $p \ge 5$ or in 3D, where edge/vertex pockets give $s$ up to
$\sqrt3$.

\subsection*{N.8 The regularity mechanism (demoted to second order)}

Reentrant staircase corners carry Kondrat'ev singularities
$r^{2/3}\sin(2\theta/3)$; they are recreated at new positions on every
level (non-nested active sets), so the associated error is $O(h)$-wavelength
by construction and falls entirely on the fine-level smoother, and
degree-$p$ approximation saturates at $O(h^{2/3})$ locally. This degrades
constants and localizes the hard error components at the boundary --- it
explains \emph{where} the smoother must work hardest, but not the observed
$p$-thresholds: the certificates of N.4 do, quantitatively and with the
right $\lambda$-, $s$-, and configuration-dependence. If the exact SBM
extrapolation were used ($k$ exact, no indefiniteness), the corner
singularities alone would be handled by the existing patch smoother, as
they are for $p \le 2$.

\subsection*{N.9 What is proven, what is certified, what is conjectured}

\emph{Proven} (Lemmas N.1--N.4): sharp Chebyshev amplification with the
per-degree factor $\rho(s)^2$; unconditional 1D $p=1$ coercivity (outward)
and the inward threshold $\sigma(h-\delta) > 1$; the kernel-escape
mechanism and penalty-insensitivity; 1D $p=2$ coercivity to $s\approx2.5$.

\emph{Certified numerically} (exact statements about finite matrices,
N.4/N.6/N.7): global indefiniteness of the symmetric part for the listed
$(p, s, \text{config})$; its penalty-insensitivity; positivity of the
chamfered $\mathbb P_3$ and truncated-$k$ variants.

\emph{Proven} (Lemma N.5): no cell-local exact extension can beat the
Chebyshev amplification; alternatives must use the PDE, the data $g$, or a
larger support (N.7b).

\emph{Conjectured / to validate in the solver}: (a) chamfer $\Rightarrow$
$p=3$ uniformly robust; (b) chamfer $+$ PDE-substituted extension
$\mathcal E_{\mathrm{pde}}$ $\Rightarrow$ $p=4$ robust \emph{at full
order} (the leading recommendation); (c) shift-adaptive $k$-truncation
alone (no chamfer) already recovers $p=4$ convergence with the existing
smoother at reduced boundary accuracy --- still the cheapest diagnostic
experiment: a ten-line change that directly tests the whole theory;
(d) restoring definiteness restores the additive patch smoother to
usefulness (with standard multiplicity damping), enabling the parallel
variant.

% --- references to add if these notes graduate into the paper:
% Rivlin, "Chebyshev Polynomials" (extremal growth outside [-1,1]).
% Grisvard, "Elliptic Problems in Nonsmooth Domains" (corner asymptotics).
% Atallah, Canuto, Scovazzi (SBM analysis, d/h conditions).
% Cai & Widlund SINUM 1992/93; Eisenstat, Elman, Schultz 1983 (nonsymmetric bounds).
% Burman et al. (CutFEM/ghost penalty, for the contrast drawn in N.6).
 at the end of sbm_patch_mg.tex to compile.
% All numerical values reproducible via theory_certificates.py (numpy only).
% ==========================================================================

\section*{Appendix (working notes): definiteness of the SBM boundary terms at the surrogate staircase, and consequences for smoothers}

\subsection*{N.0 Executive summary}

Write $s := \|\mathbf d\|/h$ for the shift length in units of the mesh size.
On a staircase surrogate ($\lambda = 0$) the flat faces have $s \lesssim 1$
and the cells at staircase steps carry \emph{two} extrapolated faces with
$s$ up to $\sqrt 2$ at reentrant pockets. The findings of these notes:

\begin{enumerate}
  \item The Taylor extension of a degree-$p$ polynomial over distance $sh$
        is polynomial extrapolation; its sharp amplification is the
        Chebyshev factor $T_p(1+2s)$, which grows by a factor
        $\rho(s)^2 = \big(\sqrt s + \sqrt{1+s}\big)^2$ \emph{per polynomial
        degree} ($\approx 3.7$ at $s=\tfrac12$, $\approx 7.5$ at
        $s=\sqrt2$). This scaling is confirmed to within a few percent by
        the computed spectra (N.4).
  \item In 1D the method is provably healthy: $p=1$ is unconditionally
        coercive for outward shifts (Lemma N.2), and $p=2$ remains coercive
        up to the absurd shift $s \approx 2.5$ (Lemma N.4). \emph{The
        breakdown is intrinsically multi-dimensional.}
  \item Computable \emph{global indefiniteness certificates} (N.4): the
        symmetric part of the SBM bilinear form is indefinite ---
        on functions supported in a single boundary cell, hence on any mesh
        containing that configuration --- for
        \[
          p = 3:\ \text{only at two-face step cells with } s \to \sqrt 2
          \quad\text{(deep reentrant pockets)},
        \]
        \[
          p = 4:\ \text{from } s \approx 0.5 \text{ at step cells, } s
          \approx 1 \text{ even at flat faces.}
        \]
        This reproduces the empirical pattern exactly: $p\le2$ robust,
        $p=3$ fails \emph{sometimes} (only when the geometry creates deep
        pockets), $p=4$ fails \emph{always} (every staircase contains
        $s \ge 0.5$ step cells). Increasing the penalty $\sigma$ does
        \textbf{not} help (N.4, item v): the offending modes nearly
        annihilate $\mathcal E$, so they do not feel the penalty.
  \item Inward shifts ($\lambda > 0$, true boundary inside the surrogate)
        are far worse: indefinite for all $p \ge 3$ at \emph{every} shift
        length tested, and not repaired by any of the fixes below. For
        $p \ge 3$ the threshold $\lambda$ should be $0$.
  \item Chamfering the staircase with triangles (the proposal of these
        notes' earlier draft) \textbf{rescues $p=3$ but not $p=4$}: it
        removes exactly the deep-pocket configurations where $p=3$ fails,
        but $p=4$ is already indefinite at flat faces with $s \approx 1$,
        which chamfering does not shorten. There is no sliver problem
        (cuts are mesh-defined, all cells uniformly shape-regular), but a
        triangle is thin in the extrapolation direction, which eats part
        of the benefit (N.6).
  \item Truncating the Taylor order to $k=2$ ($k=1$) restores definiteness
        at step cells up to $s \approx 0.75$ (everywhere), at the price of
        reduced boundary consistency $O(h^{k+1})$ --- which defeats the
        purpose of $p=4$ (N.7).
  \item A no-go lemma (N.7b): \emph{any} linear cell-local extension exact
        on $\mathbb Q_p(K)$ coincides with polynomial extrapolation and
        carries the full Chebyshev amplification. Alternatives must import
        extra information: the PDE, the boundary data $g$, or a larger
        support. The PDE-substituted extension (all normal derivatives of
        order $\ge 2$ traded for tangential derivatives and $f$-data)
        keeps full consistency order, gains roughly one polynomial degree
        of definiteness (certified in N.7b), and combined with the chamfer
        (which caps $s$ at $\approx 0.7$ and removes two-face cells) is
        the predicted robust \emph{full-order} $p=4$ configuration.
\end{enumerate}

\subsection*{N.1 Setting}

Model: cell $K = (0,h)^d$, $\mathbb Q_p$ Lagrange polynomials, SBM form
restricted to $K$ and its surrogate faces $F$:
\begin{equation}
  a_K(u,v) = \int_K \nabla u\cdot\nabla v
  - \int_F (\partial_n u)\, v
  - \int_F (\partial_n v)\, \mathcal E u
  + \sigma \int_F \mathcal E u\, \mathcal E v,
  \qquad \sigma = \gamma\, p^2/h,
  \label{eq:cellform}
\end{equation}
with $\mathcal E u(\mathbf x) = \sum_{j=0}^{k} \tfrac1{j!} (\mathbf
d\cdot\nabla)^j u(\mathbf x)$. Throughout, $s = \|\mathbf d\|/h$ and we set
$h=1$ in computations; $\gamma = 4$ unless stated (results are insensitive
to $\gamma$, see N.4).

Two mechanisms could plausibly cause the observed breakdown: (i) elliptic
corner singularities of the surrogate domain (reentrant staircase corners,
interior angle $3\pi/2$, Kondrat'ev exponent $r^{2/3}$), and (ii) loss of
definiteness of the discrete boundary terms. The results below show that
(ii) is dominant and quantitatively predictive; (i) is discussed briefly in
N.8.

\subsection*{N.2 Sharp extrapolation amplification}

\paragraph{Lemma N.1 (extension = extrapolation; sharp growth).}
\emph{Let $q$ be a polynomial of degree $p$ on $[0,h]$ and let the shift be
axis-aligned. For $k \ge p$ the Taylor extension is exact point evaluation,
$\mathcal E q(0) = q(-\delta)$, and}
\[
  \sup\Big\{ |q(-\delta)| \,:\, \|q\|_{L^\infty(0,h)} \le 1 \Big\}
  \;=\; T_p\!\big(1 + 2s\big), \qquad s = \delta/h,
\]
\emph{attained by the Chebyshev polynomial. Moreover, with
$\rho(s) := \sqrt s + \sqrt{1+s}$,}
\[
  T_p(1+2s) \;=\; \tfrac12\Big( \rho(s)^{2p} + \rho(s)^{-2p} \Big).
\]

\emph{Proof.} Exactness: the Taylor polynomial of order $k\ge p$ of a
degree-$p$ polynomial is the polynomial itself. Extremality of $T_p$
outside $[-1,1]$ is classical (Chebyshev; see Rivlin). The identity follows
from $T_p(\cosh\theta) = \cosh(p\theta)$ with
$1+2s = \cosh\big(2\,\mathrm{arcsinh}\sqrt s\big)$ and
$e^{\mathrm{arcsinh}\sqrt s} = \sqrt s + \sqrt{1+s}$. $\square$

The quantity $\rho(s)^2$ is the \emph{amplification per polynomial degree}:

\begin{center}
\begin{tabular}{lccccc}
\hline
 & $s=\tfrac14$ & $s=\tfrac12$ & $s=\tfrac34$ & $s=1$ & $s=\sqrt2$ \\
\hline
$\rho(s)^2$ (per degree)      & $2.6$ & $3.7$ & $4.8$ & $5.8$ & $7.5$ \\
$T_1(1+2s)$ & $1.5$ & $2$ & $2.5$ & $3$ & $3.8$ \\
$T_2(1+2s)$ & $3.5$ & $7$ & $11.5$ & $17$ & $28$ \\
$T_3(1+2s)$ & $9$ & $26$ & $55$ & $99$ & $214$ \\
$T_4(1+2s)$ & $24.5$ & $97$ & $265$ & $577$ & $1611$ \\
\hline
\end{tabular}
\end{center}

The penalty and consistency terms in \eqref{eq:cellform} are bilinear in
$\mathcal E$, so the boundary block scales like $\sigma\, T_p^2$ relative to
the $O(1)$ (in $h^{d-2}$ units) volume stiffness. This is confirmed by the
computed largest eigenvalues $\lambda_{\max}$ of the symmetric part: e.g.\
for the flat configuration at $p=4$, the measured ratios
$\lambda_{\max}(s{=}1)/\lambda_{\max}(s{=}\tfrac12) = 34$ and
$\lambda_{\max}(s{=}\sqrt2)/\lambda_{\max}(s{=}1) = 7.6$ match the predicted
$\big(T_4(3)/T_4(2)\big)^2 = 35$ and $\big(T_4(3.83)/T_4(3)\big)^2 = 7.8$.
At $p=4$, $s=1$ the local block condition number already exceeds $10^7$.
Conditioning alone, however, is not the failure mode --- definiteness is,
as the next two subsections show.

\subsection*{N.3 Exact one-dimensional results}

On $K=(0,h)$ with surrogate face at $x=0$ and true boundary at $-\delta$
(outward shift), $\mathcal E v(0) = v(-\delta)$, form \eqref{eq:cellform}.

\paragraph{Lemma N.2 ($p=1$, closed form).} \emph{For $v = a + bx$:}
\[
  a_K(v,v) = b^2 (h+\delta) + 2bw + \sigma w^2, \qquad w := v(-\delta),
\]
\emph{which is positive definite iff $\sigma (h + \delta) > 1$. For the
inward shift ($\Gamma$ at $+\delta$ inside the cell) the same computation
gives $b^2(h-\delta) + 2bw + \sigma w^2$, positive definite iff
$\sigma (h - \delta) > 1$.}

Consequences: $p=1$ with outward shifts is unconditionally coercive (any
$\sigma \ge 1/h$, any $\delta > 0$) --- matching the total robustness of
$p=1$ in practice. Inward shifts lose coercivity as $\delta \to h$; with
$\sigma = 4/h$ the threshold is $\delta = 3h/4$, reproduced to machine
precision by the 2D code (a useful cross-validation of both).

\paragraph{Lemma N.3 (kernel escape).} \emph{On the subspace
$\{v : \mathcal E v = 0 \text{ on } F\}$ the penalty and one flux term
vanish identically and \eqref{eq:cellform} reduces to the
\emph{unpenalized half-Nitsche} form}
\[
  F(v) = \|\nabla v\|_{L^2(K)}^2 - \int_F (\partial_n v)\, v .
\]
Definiteness on this kernel is a sign-correlation property between $v$ and
$\partial_n v$ on $F$, forced by $v$ vanishing at the shifted points; the
penalty is powerless there, no matter how large. This is the structural
reason for the $\gamma$-insensitivity observed in N.4.

\paragraph{Lemma N.4 ($p=2$, 1D, explicit threshold).} \emph{For $v$
quadratic with $v(-\delta) = 0$ (kernel of $\mathcal E$), writing
$c = v'(0) $-parametrization $v = (x+\delta)(a+bx)$, $h=1$:}
\[
  F(v) = (1+s)\,c^2 + (2 - s^2)\,cb + \tfrac43\, b^2, \qquad c = a + bs,
\]
\emph{positive definite iff $(2-s^2)^2 < \tfrac{16}{3}(1+s)$, i.e.\ for all
$s \lesssim 2.5$.}

So in 1D even $p=2$ survives shifts of two and a half cells; the observed
2D breakdown at $s \le \sqrt2$ cannot be a 1D phenomenon. What 1D misses:
at a staircase step, \emph{two} faces with orthogonal shift directions load
the same cell polynomial, and the sign-correlation of Lemma N.3 cannot be
satisfied for both directions simultaneously once $p$ gives the polynomial
enough freedom to oscillate between the cell and the shifted evaluation
points.

\subsection*{N.4 Numerical definiteness certificates in 2D}

\paragraph{Certificate construction.} Let $K$ be a boundary cell in one of
the configurations below and let $v \in \mathbb Q_p(K)$ vanish on the edges
of $K$ shared with active neighbors. Extending $v$ by zero yields a
conforming global function whose global energy equals $a_K(v,v)$: volume
terms of neighbors vanish, and all other surrogate-face terms vanish
because $\mathcal E$ on a face of cell $K'$ uses only the polynomial of
$K'$. Hence
\[
  \lambda_{\min}\Big( \tfrac12 (A_K + A_K^T)\big|_{S} \Big) < 0
  \;\Longrightarrow\;
  \text{the \emph{global} symmetric part is indefinite}
\]
on every mesh containing the configuration, where $S$ is the zero-extension
subspace. Configurations ($h=1$, exact extrapolation $=$ Taylor with $k=p$
for the axis-aligned shifts):
\emph{flat} (one face $\{y{=}1\}$, $\mathbf d = (0,s)$; $S$: $v=0$ on the
three interior edges), \emph{step} (two faces $\{y{=}1\},\{x{=}1\}$,
$\mathbf d = (0,s),(s,0)$; the staircase step cell; $S$: $v=0$ on two
interior edges), \emph{step-diag} (as step but both shifts
$s(1,1)/\sqrt2$, modeling pocket-directed closest-point projections), and
\emph{chamfer} variants in N.6.

\paragraph{Results ($\gamma = 4$).} Minimal eigenvalue of the symmetric
part on $S$ (negative $=$ certified global indefiniteness; the volume
stiffness is $O(1)$ in these units, so magnitudes are meaningful):

\begin{center}
\begin{tabular}{llcccc}
\hline
config & $p$ & $s=0.5$ & $s=0.75$ & $s=1$ & $s=\sqrt2$ \\
\hline
flat  & 3 & $+1.22$ & $+0.91$ & $+0.75$ & $+0.61$ \\
flat  & 4 & $+0.28$ & $+0.08$ & $\mathbf{-0.03}$ & $\mathbf{-0.13}$ \\
step  & 3 & $+0.27$ & $+0.09$ & $+0.02$ & $\mathbf{-0.05}$ \\
step  & 4 & $\mathbf{-0.04}$ & $\mathbf{-0.13}$ & $\mathbf{-0.18}$ & $\mathbf{-0.24}$ \\
\hline
\end{tabular}
\end{center}

Findings:
\begin{enumerate}
  \item[(i)] $p \le 2$: positive in every configuration and shift tested
        (not shown; margins $O(1)$).
  \item[(ii)] $p=3$: indefinite \emph{only} at the two-face step cell with
        $s \to \sqrt2$, i.e.\ at deep reentrant pockets. Whether a given
        geometry produces such pockets depends on how $\Gamma$ cuts the
        grid --- hence ``$p=3$ sometimes breaks''.
  \item[(iii)] $p=4$: indefinite from $s \approx 0.5$ at step cells and
        from $s \approx 1$ at flat faces. Every staircase surrogate of a
        generic $\Gamma$ contains such cells --- hence ``$p=4$ always
        breaks''.
  \item[(iv)] The nonsymmetric spectra show only small imaginary parts for
        outward shifts: for $\lambda = 0$ the pathology is
        \emph{indefiniteness of the symmetric part}, not complex
        eigenvalues (those belong to inward shifts, consistent with the 1D
        study in~\cite{wichrowski2025geometric}).
  \item[(v)] Penalty insensitivity (Lemma N.3 in action), step cell,
        $s=1$: for $p=4$, $\lambda_{\min}|_S = -0.191, -0.184, -0.178,
        -0.177$ at $\gamma = 2, 4, 16, 64$. No penalty tuning can repair
        this.
  \item[(vi)] Inward shifts (flat face, $\mathbf d = (0,-s)$, corrected
        certificate): $p=2$ positive at all $s$; $p=3$ and $p=4$
        \emph{negative at every} $s \in \{0.25, 0.5, 0.75\}$ (e.g.\ $-0.63$
        and $-1.13$ at $s=0.25$). Intersected-cell surrogates
        ($\lambda>0$) are therefore incompatible with $p \ge 3$
        independently of corner effects.
\end{enumerate}

\subsection*{N.5 Consequences for smoothers}

Let $v_\star$ be a certified negative mode ($a(v_\star,v_\star) < 0$,
supported in one boundary cell).

\begin{itemize}
  \item \textbf{All energy-based convergence theory is void.} Both the
        Schwarz theory for additive/multiplicative smoothers and the
        field-of-values GMRES bounds (Eisenstat--Elman--Schultz,
        Cai--Widlund) presuppose a definite symmetric part. On staircase
        SBM with $p \ge 3$ this hypothesis \emph{fails as a matter of
        fact}, not of proof technique.
  \item \textbf{Pointwise and cell-wise additive smoothers.} For
        (block-)Jacobi-type methods the local blocks themselves become
        indefinite (the single-cell Neumann blocks go negative one degree
        earlier than the certificates, already $p=3$, $s \ge 0.75$).
        Damped Richardson with an indefinite block-diagonal preconditioner
        has no admissible damping: the modes with negative Rayleigh
        quotient are amplified for every $\omega > 0$ of either sign
        globally. This is the cleanest available explanation for the
        observed \emph{divergence} (not merely stagnation) of the
        cell-wise additive DG+SBM smoother.
  \item \textbf{Why multiplicative full-residual patches still work at
        $p \le 3$.} The negative mode $v_\star$ is supported in one cell,
        hence entirely inside at least one vertex patch (the full-residual
        construction guarantees its whole residual row is visible there).
        An \emph{exact} local solve needs only invertibility, not
        definiteness, of $A_i$; visiting that patch removes $v_\star$'s
        error component exactly, and the sequential sweep propagates the
        repaired residual outward. In the additive version the $2^d$
        overlapping patches containing $v_\star$ each compute a correction
        against an indefinite, non-normal block and their \emph{sum} is
        applied blindly; there is no self-correction within the sweep, and
        the usual multiplicity damping $\omega \le 1/N_c$ cannot produce
        contraction where the field of values crosses zero.
  \item \textbf{Why $p=4$ defeats even the multiplicative smoother.} The
        indefinite modes are no longer isolated: at $p=4$ they exist at
        \emph{every} boundary cell with $s \gtrsim 0.5$--$1$, they couple
        to neighboring patches through off-block entries of size
        $O(\sigma T_p^2)$ ($\ge 10^6$ at $s = 1$), and the error injected
        across a patch boundary is re-amplified faster than the sweep
        damps it. Locality of the repair mechanism is lost.
\end{itemize}

\subsection*{N.6 Chamfer (mesh-defined triangular cuts): rescues $p=3$, not $p=4$}

The chamfer splits staircase corner quads along the vertex-to-vertex
diagonal. The cut position is fixed by the mesh: every cut cell is exactly
half a quad, uniformly shape-regular --- \emph{there is no small-cut/sliver
problem and no ghost penalty is needed}, in contrast to
$\Gamma$-interpolating CutFEM cuts. Reentrant $3\pi/2$ corners become pairs
of $5\pi/4$ corners (regularity exponent $2/3 \to 4/5$), and the worst
shift drops from $s \approx \sqrt2$ to $s \approx 0.7$ with a single,
near-normal face replacing the two-face loading.

Certificates for the chamfered cell (hypotenuse face, shift
$s(1,1)/\sqrt2$), two variants: $\mathbb P_p$ on the triangle (remeshed)
and $\mathbb Q_p$ on the parent quad integrated over the triangle
(cut-quadrature):

\begin{center}
\begin{tabular}{llccc}
\hline
variant & $p$ & $s=0.5$ & $s=0.707$ & $s=1$ \\
\hline
$\mathbb P_p$ triangle & 3 & $+0.91$ & $+0.62$ & $+0.42$ \\
$\mathbb P_p$ triangle & 4 & $\mathbf{-0.15}$ & $\mathbf{-0.34}$ & $\mathbf{-0.48}$ \\
$\mathbb Q_p$ cut      & 3 & $\mathbf{-0.01}$ & $\approx\mathbf{-0.03}$ & $\mathbf{-0.02}$ \\
$\mathbb Q_p$ cut      & 4 & $\mathbf{-0.08}$ & $\mathbf{-0.07}$ & $\mathbf{-0.06}$ \\
\hline
\end{tabular}
\end{center}

Reading:
\begin{itemize}
  \item The \emph{remeshed} chamfer ($\mathbb P_p$ triangles) is positive
        for $p=3$ at all relevant shifts. Since the chamfer also caps $s$
        at $\approx 0.7$, it removes exactly the configuration in which
        $p=3$ was certified indefinite. \textbf{Prediction: chamfering
        makes $p=3$ uniformly robust.}
  \item It does \emph{not} rescue $p=4$: the triangle block is indefinite
        from $s = 0.5$. Two reasons: (a) a triangle is thin in the
        extrapolation direction (extent $h/\sqrt2$ normal to the
        hypotenuse, tapering), so the effective Chebyshev argument is
        larger than $s$ suggests; (b) independently, $p=4$ is already
        indefinite at \emph{flat} faces with $s \approx 1$, which the
        chamfer does not touch.
  \item The cut-quadrature variant is weaker (marginally indefinite even
        at $p=3$): the polynomial is energy-controlled only on the
        half-cell while being extrapolated from its edge. If chamfering is
        implemented, the triangles should carry genuine $\mathbb P_p$
        spaces (or the parent-cell variant needs a stabilization on the
        inactive half). This is an implementation-relevant distinction the
        earlier draft of these notes missed.
\end{itemize}

\subsection*{N.7 The effective fix: truncated extrapolation order}

Lemma N.1 identifies the amplification driver as the extrapolation
\emph{order}: truncating the Taylor extension at order $k < p$ caps the
growth at degree-$k$ extrapolation ($T_k$-like) while retaining
$O((sh)^{k+1})$ boundary consistency. Certificates for the step cell
(worst outward configuration), truncation order $k$:

\begin{center}
\begin{tabular}{llcccc}
\hline
config & $p$ & $k=p$ & $k=3$ & $k=2$ & $k=1$ \\
\hline
step, $s=0.75$   & 4 & $-0.13$ & $-0.10$ & $\mathbf{+0.01}$ & $+0.41$ \\
step, $s=1$      & 4 & $-0.18$ & $-0.14$ & $-0.03$ & $+0.35$ \\
step, $s=\sqrt2$ & 4 & $-0.24$ & $-0.18$ & $-0.06$ & $+0.30$ \\
step, $s=\sqrt2$ & 3 & $-0.05$ & --      & $\mathbf{+0.04}$ & $+0.50$ \\
flat, $s=1$      & 4 & $-0.03$ & $+0.06$ & $+0.34$ & $+0.98$ \\
\hline
\end{tabular}
\end{center}

\begin{itemize}
  \item $k=1$ restores definiteness \emph{everywhere} with $O(1)$ margins
        --- a bulletproof but low-order option (locally $O(h^2)$
        consistency; acceptable if applied only at the few worst faces).
  \item $k=2$ restores definiteness for $p=4$ up to $s \approx 0.75$ and
        for $p=3$ everywhere, with $O(h^3)$ boundary consistency.
  \item Honest objection: with a staircase or chamfered surrogate,
        $\delta = O(h)$, so $k<p$ caps the \emph{global} boundary accuracy
        at $O(h^{k+1})$ and largely defeats the purpose of $p=4$. (The
        truncation is order-neutral only where $\delta = O(h^2)$, i.e.\
        with a $\Gamma$-interpolating cut.) A shift-adaptive rule
        ($k=p$ where $s \le 0.4$, smaller $k$ beyond) confines the loss to
        the worst faces, but the principled resolution is to change the
        extension, not its order --- see N.7b.
  \item For inward shifts truncation helps at moderate $s$ but $p \ge 3$
        remains indefinite at small inward $s$ regardless of $k$ ---
        reinforcing finding (vi): use $\lambda = 0$ for $p \ge 3$.
\end{itemize}

\subsection*{N.7b Alternative extension operators: a no-go lemma and three ways around it}

\paragraph{Lemma N.5 (no free lunch for cell-local extensions).}
\emph{Let $\mathcal L : \mathbb Q_p(K) \to \mathbb R$ be linear, computed
from data on $K$ only, and exact on polynomials,
$\mathcal L q = q(\tilde{\mathbf x} + \mathbf d)$ for all $q \in
\mathbb Q_p(K)$ (as required for $O(h^{p+1})$ consistency from cell-local
data). Then $\mathcal L$ \emph{is} the point-evaluation extrapolation on
$\mathbb Q_p(K)$, and its norm is the Chebyshev factor of Lemma N.1.}
\emph{Proof:} two linear functionals agreeing on the whole space are
equal. $\square$

Smoothness of the extension is a red herring: on the discrete space every
exact extension coincides with the one already in use. Any genuine
alternative must therefore use information beyond the single cell's
polynomial. Exactly three sources exist, and each gives a known method
family:

\paragraph{(A) Use the PDE: substituted normal derivatives (recommended).}
Since $-\Delta u = f$, all normal derivatives of order $\ge 2$ can be
eliminated recursively,
\[
  \partial_n^2 u = -f - \Delta_{\mathrm{tan}} u, \qquad
  \partial_n^3 u = -\partial_n f - \Delta_{\mathrm{tan}} \partial_n u,
  \quad \dots
\]
so the degree-$k$ Taylor extension becomes
\[
  \mathcal E_{\mathrm{pde}} u
  = u + \delta\, \partial_n u
  + \sum_{j=2}^{k} \tfrac{\delta^j}{j!}
    \big( \text{tangential derivatives of } u,\ \partial_n u \big)
  + \underbrace{\big( f\text{-terms} \big)}_{\text{data, rhs}} ,
\]
with \emph{every $u$-dependent term evaluated on the face}: no
extrapolation beyond the first normal derivative. Crucially,
$\mathcal E_{\mathrm{pde}}$ is exact on solutions of the PDE --- which is
all that consistency requires --- so \textbf{the full order $O(\delta^{k+1})$
is retained} while the Chebyshev growth is traded for tangential Markov
factors $\big(c\,p^2 \delta/h\big)^{j}$, small for moderate $s$.
Certificates (same setup as N.4, $f$-part irrelevant for the spectrum):

\begin{center}
\begin{tabular}{llcccc}
\hline
config & $p$ & $s=0.5$ & $s=0.75$ & $s=1$ & $s=\sqrt2$ \\
\hline
flat, $\mathcal E_{\mathrm{pde}}$ & 3 & $+2.87$ & $+3.36$ & $+3.84$ & $+4.25$ \\
flat, $\mathcal E_{\mathrm{pde}}$ & 4 & $+1.34$ & $+0.80$ & $-0.01$ & $-0.71$ \\
step, $\mathcal E_{\mathrm{pde}}$ & 3 & $+1.20$ & $+0.90$ & $+0.49$ & $-0.10$ \\
step, $\mathcal E_{\mathrm{pde}}$ & 4 & $+0.12$ & $-0.34$ & $-0.38$ & $-0.33$ \\
\hline
\end{tabular}
\end{center}

Compared with plain Taylor (N.4): roughly one polynomial degree of
definiteness is gained at full consistency order ($p=4$ flat: $+0.08 \to
+0.80$ at $s = 0.75$; step: indefinite from $s=0.5 \to$ definite at
$s=0.5$). At $s \gtrsim 1$ the tangential factors take over and
indefiniteness returns --- so $\mathcal E_{\mathrm{pde}}$ alone does not
cure the staircase, but \textbf{chamfer ($s \le 0.7$, single near-normal
face) $+$ $\mathcal E_{\mathrm{pde}}$ is certified-definite in every
configuration that then occurs, at full order.} Implementation cost:
tangential/mixed derivatives of the shape functions on the face (already
available), plus $f$ and its normal-derivative data near the boundary in
the right-hand side.

\paragraph{(B) Use the boundary data: interpolatory (ghost/ROD-type)
conditions.} Taylor never uses the one thing that is known exactly: $g$ at
the projected point $\mathbf x \in \Gamma$. Writing the boundary condition
via a 1D interpolant along the shift line through $p{+}1$ points
\emph{inside} $\tilde\Omega$ plus the pinned value $g(\mathbf x)$, the
surrogate-face value $u(\tilde{\mathbf x})$ is expressed by
\emph{interpolation} (the evaluation point lies between the samples and
$\mathbf x$), so the $u$-dependent weights are bounded by a Lebesgue
constant instead of $T_p$; full order is retained. This is the
finite-element analogue of ghost-point methods (Gibou--Fedkiw) and
reconstructed off-site data (ROD). Cost: the stencil crosses 1--2 interior
cells, so the sparsity pattern widens; the samples stay inside a vertex
patch, so the full-residual construction survives.

\paragraph{(C) Use a larger support: aggregated extension.} Extrapolate
from a least-squares polynomial fit over an interior patch of diameter
$L = 2h$--$3h$ (AgFEM-flavored). The amplification drops to
$T_p(1 + 2\delta/L)$; by the certificate tables, an effective shift
$\delta/L \le 0.4$ is safe even at $p=4$, so $L \ge 3h$ suffices for the
worst staircase pockets. Cost: same widened coupling as (B), plus the fit;
conceptually the bluntest of the three.

Ranking for this project: (A) is the least invasive (face-local, no
sparsity change, no new data structures) and combines multiplicatively
with the chamfer; (B) is the most robust in principle (bounded weights
regardless of geometry) and is the natural fallback if (A)$+$chamfer still
misbehaves at $p \ge 5$ or in 3D, where edge/vertex pockets give $s$ up to
$\sqrt3$.

\subsection*{N.8 The regularity mechanism (demoted to second order)}

Reentrant staircase corners carry Kondrat'ev singularities
$r^{2/3}\sin(2\theta/3)$; they are recreated at new positions on every
level (non-nested active sets), so the associated error is $O(h)$-wavelength
by construction and falls entirely on the fine-level smoother, and
degree-$p$ approximation saturates at $O(h^{2/3})$ locally. This degrades
constants and localizes the hard error components at the boundary --- it
explains \emph{where} the smoother must work hardest, but not the observed
$p$-thresholds: the certificates of N.4 do, quantitatively and with the
right $\lambda$-, $s$-, and configuration-dependence. If the exact SBM
extrapolation were used ($k$ exact, no indefiniteness), the corner
singularities alone would be handled by the existing patch smoother, as
they are for $p \le 2$.

\subsection*{N.9 What is proven, what is certified, what is conjectured}

\emph{Proven} (Lemmas N.1--N.4): sharp Chebyshev amplification with the
per-degree factor $\rho(s)^2$; unconditional 1D $p=1$ coercivity (outward)
and the inward threshold $\sigma(h-\delta) > 1$; the kernel-escape
mechanism and penalty-insensitivity; 1D $p=2$ coercivity to $s\approx2.5$.

\emph{Certified numerically} (exact statements about finite matrices,
N.4/N.6/N.7): global indefiniteness of the symmetric part for the listed
$(p, s, \text{config})$; its penalty-insensitivity; positivity of the
chamfered $\mathbb P_3$ and truncated-$k$ variants.

\emph{Proven} (Lemma N.5): no cell-local exact extension can beat the
Chebyshev amplification; alternatives must use the PDE, the data $g$, or a
larger support (N.7b).

\emph{Conjectured / to validate in the solver}: (a) chamfer $\Rightarrow$
$p=3$ uniformly robust; (b) chamfer $+$ PDE-substituted extension
$\mathcal E_{\mathrm{pde}}$ $\Rightarrow$ $p=4$ robust \emph{at full
order} (the leading recommendation); (c) shift-adaptive $k$-truncation
alone (no chamfer) already recovers $p=4$ convergence with the existing
smoother at reduced boundary accuracy --- still the cheapest diagnostic
experiment: a ten-line change that directly tests the whole theory;
(d) restoring definiteness restores the additive patch smoother to
usefulness (with standard multiplicity damping), enabling the parallel
variant.

% --- references to add if these notes graduate into the paper:
% Rivlin, "Chebyshev Polynomials" (extremal growth outside [-1,1]).
% Grisvard, "Elliptic Problems in Nonsmooth Domains" (corner asymptotics).
% Atallah, Canuto, Scovazzi (SBM analysis, d/h conditions).
% Cai & Widlund SINUM 1992/93; Eisenstat, Elman, Schultz 1983 (nonsymmetric bounds).
% Burman et al. (CutFEM/ghost penalty, for the contrast drawn in N.6).
